% The appendix of the simulated readers: the texts sent to the models, the probes and the measures the body does not print. The texts are copied from src/mootcourt/simulators/readers/protocol.py, which was fixed before the first call; the tables come from paper/summary_tables/readers.py.
\section{Simulated Readers in Detail}
\label{sec:readers-detail}
% src: the numbers 1 to 5 are the numbers of the methods in the list of this appendix; no measured number
This appendix gives what Section~\ref{sec:readers} leaves out: how a feed is built, the panel, the references and the seven methods in full; the two runs and how they compare; the run with an open model, its direct judgment and the embedding classifier; the readings with the quarters as the units; the seven methods side by side; the three probes for earlier exposure to the archive; the texts sent to the models; how methods 1 to 5 are fitted and how the calls were made; and the measures that carry no verdict.

% src: paper/data/run11-simulated-readers/plan.json counts.tests 1770, counts.contrasts 5966, counts.blocks 63, counts.feeds 308, block_size 30, readers 10; tables/fmt-readers-more.tex (units ranked 1,771)
\paragraph{One reader, one feed.} The rows are the main run's scored rows in the headline-only cell, where the arms of a test differ in the headline alone: 1,770 tests with 5,966 contrasts. In the real test a visitor saw one package and clicked or did not. A simulated reader therefore never sees two headlines of one test. One call to the model plays one reader who scrolls one feed of up to 30 story cards. Each card shows the headline of one arm, each from a different test, and nothing else. For each card the reader returns a probability from 0 to 100 that they would click it. The tests of a scored origin are cut into blocks of 30. A story is a test, and a block has as many feeds as its largest test has arms. Each feed of a block shows one arm of every test of the block, and each test keeps the same position in every feed. The arms of a test are spread over the feeds, one per feed. Two arms of a test are thus seen in the same place among the same neighbouring stories and differ only in the headline. A test with fewer arms repeats them to fill the feeds, and an arm's score averages its repeats. This gives 63 blocks and 308 feeds.

% src: paper/data/run13-readers-open-model/README.md (run 13 is the name of its run folder); paper/data/run11-simulated-readers/README.md (the three models, their training data cutoff of June 2026, the primary model declared before the run); src/mootcourt/simulators/readers/protocol.py SYSTEM_PROMPT (between 2013 and 2015; no publisher named), PROFILES (ten), HALVES; no measured number; 5.5 is part of each model's name
\paragraph{The panel.} Ten reader profiles read every feed. A profile is a short paragraph that gives an age band, an occupation, interests, a mood, a device and how the reader met the story. The archive holds no audience data, so the profiles are assumed, and their spread is variation between prompts and not a sample of the real audience. The instructions place the reader between 2013 and 2015 and do not name the publisher. Each call's probabilities are centred on that call's own mean. Centring removes the general level of a reader and of a feed, so that only differences between headlines remain. An arm's \emph{panel score} is the mean of the centred probabilities over the readers. For a contrast, the panel's direction is the sign of the variant's score minus the reference's, and a difference of exactly zero counts as half right. We ran three models of one vendor, from the largest to the smallest Claude Opus 5.5, Claude Sonnet 5.5 and Claude Haiku 5.5, and declared the largest as primary before the run. A fourth model, open and from another maker, was added after the results (run 13) and is not primary. The text calls the three the primary or largest model, the middle model and the smallest model, and the tables call them Opus, Sonnet and Haiku. The archive was public years before the training data of all three ends, so nothing is hidden from them, and the models may have met these headlines. Three probes, described below, describe that risk and cannot exclude it.

% src: tables/fmt-readers-readings.tex rows the earliest-arm rule 0.5701, the text classifier 0.5484, the direction ceiling 0.8301; paper/data/run11-simulated-readers/simulated-readers-claude-opus-5-5-20261106-2dc461f52ea4.json baselines.always_reference.accuracy 0.5701, baselines.text_lr.accuracy 0.5484, baselines.ceiling.accuracy 0.8301; tables/fmt-direction.tex row always_reference, accuracy 0.582; paper/data/run9-visible-exposure/backtest-upworthy-20261106-2dc461f52ea4.json tables.aggregate[cell=ALL,rung=always_reference].direction_accuracy 0.582 (all three cells), tables.aggregate[cell=headline_only,rung=always_reference].direction_accuracy 0.570
\paragraph{References.} The measure is \emph{direction accuracy}: the share of contrasts on which a method names the sign of the observed effect. Three references are computed on the same rows. The \emph{earliest-arm rule} always says that the variant loses to the reference arm, the package created first, and it is right on 0.570 of the contrasts. It is the baseline that Section~\ref{sec:method} calls always the reference arm. On the contrasts of all three cells the same rule is right on 0.582 (Table~\ref{tab:fmt-direction}), and the figure here is lower because it covers the headline-only cell alone. The text classifier of Section~\ref{sec:method}, fitted on earlier tests, is right on 0.548. The ceiling of Section~\ref{sec:protocol}, an estimate of the accuracy a method would reach if it knew each test's true effect, is 0.830. A rule that always picks the first headline written thus scores 0.570 without reading anything, so one half is not the reference to beat.

% src: docs/superpowers/specs/2026-11-18-llm-simulated-readers-design.md, Aligning the readers with earlier results (seven methods; refitted at every scored origin on the headline-only tests that the visibility rule lets a forecast see and for which panel answers exist, which are the scored tests of earlier origins past their assigned end; ten earlier tests in the prompt; 200 earlier tests and a note of at most 150 words); src/mootcourt/simulators/readers/align.py visible_before (rows of an earlier scored origin whose result was usable by the origin); src/mootcourt/simulators/readers/incontext.py example_units, brief_units (the same rule; one draw per origin; every visible unit when fewer than 200 exist); src/mootcourt/simulators/readers/textfeature.py (method 5 is fitted on every visible row of every cell; a training row has a panel difference only if it was a scored headline-only row of an earlier origin, and any other row gets zero); src/mootcourt/simulators/readers/measures.py _with_plain_at_first_origin; paper/data/run11-simulated-readers/plan.json counts.by_origin.2013-04-01.earlier_units_visible 0; src/mootcourt/simulators/readers/protocol.py N_EXAMPLES 10, N_BRIEF_UNITS 200, NOTE_WORDS 150; the numbers 1 to 7 are the methods' numbers in the list below
\paragraph{Seven ways to align the panel.} Each method is refitted at every scored origin. It may use only headline-only tests that were scored at an earlier origin and whose results are visible at the present origin under the rule of Section~\ref{sec:setting}. This is a smaller history than the ladder and the text classifier of Section~\ref{sec:method} read. Tests from before the first scored origin are visible, but they were not sent to the models and have no panel answers. Methods 1 to 4 do not use them, and method 5 keeps them only as training rows of its classifier, with no panel value. Methods 6 and 7 need no panel answers and are held to the same history all the same. At the first scored origin no earlier scored test exists. Methods 1 to 4, 6 and 7 therefore equal the plain panel there. Method 5 equals the text classifier alone there, so the third reading below is zero by construction on the contrasts of that origin. No prompt, profile or setting was tuned on an outcome. Five methods use the stored answers of the plain panel.
\begin{enumerate}
\item \emph{Direction map.} A logistic regression of the observed direction on the panel's difference.
\item \emph{Reader weights.} Weights over the ten readers, fitted so that the weighted panel explains the click counts of earlier arms.
\item \emph{Audience matching.} Weights chosen only so that the panel's mean level matches the mean click rate of earlier arms, with no use of which arm won.
\item \emph{Error correction.} A monotone map from panel score to click rate, then a linear correction on four fixed traits of a headline.
\item \emph{Text classifier with the panel.} The panel's difference enters the text classifier as one more feature. A training row carries a panel difference only if it was a scored headline-only row of an earlier origin, and every other training row takes the value zero for that feature.
\end{enumerate}
Two methods put earlier results in front of the reader and need new calls, on the feeds of the plain panel.
\begin{enumerate}\setcounter{enumi}{5}
\item \emph{Earlier results in the prompt.} Before the feed, the reader sees ten earlier tests with the click rate of each headline, the same ten for every feed and reader of an origin.
\item \emph{Audience brief.} The model reads up to 200 earlier tests with their click rates and writes a note of at most 150 words on what this audience clicked. Every reader of that origin receives the note.
\end{enumerate}

% src: paper/data/run11-simulated-readers/README.md (the transport; what the tool adds: the present date beside the instruction to play a person of 2013 to 2015; no control of sampling); paper/data/run12-simulated-readers-api/README.md (the same 9,007 requests per model through the vendor's API; four request fields; no sampling setting can be sent; planned before any reply of run 11 was scored; The code: the commit of run 12 adds the second transport to the code state of run 11, which had been scored before any of it was written; Deviations from the plan: none); docs/superpowers/plans/2026-11-19-simulated-readers-api-run.md, The comparison between the two runs, fixed now; tables/fmt-readers-more.tex row planned 9,007; run 11 and run 12 are the names of the two run folders
\paragraph{Two runs of the same requests.} The first run (run 11) reached the models through the vendor's command line tool. That tool puts its own context around each call, the present date included, so a reader was told the year 2026 beside the instruction to play a person of 2013 to 2015. We therefore planned a second run (run 12) before any reply of the first was scored. It repeats the same 9,007 planned calls per model through the vendor's programming interface, where the model receives our two texts and nothing else. The calls of run 12 were made after run 11 had been scored, and no method, prompt, seed or rule changed in between. The reader calls of method 7 carry the note written in that run, so their text differs between the runs, and every other request is identical. Neither route lets us fix the sampling of a reply, so two runs of one request differ by chance as well. The rule for reading the second run was fixed with its plan: if, for the primary model, the interval of the difference in the plain panel's accuracy between the runs contains zero, we say that the added context did not move that accuracy by an amount this design can detect. The fixed words of the three readings come from run 11 in either case.

\begin{table*}[t]\centering\footnotesize
% src: tables/fmt-readers-methods.tex; paper/data/run11-simulated-readers/simulated-readers-claude-opus-5-5-20261106-2dc461f52ea4.json methods.audience_brief.minus_always_reference.resamples 500, counts.blocks 63; run 11 and run 12 are the names of the two run folders
\caption{The plain panel and the seven alignment methods in run 11, by model. The plain panel is the mean of the ten readers as they answered. Methods 1 to 5 are refitted on earlier tests from the plain panel's stored answers, and methods 6 and 7 put earlier results in the reader's prompt; the paragraph on the seven ways to align the panel describes each. The first column is direction accuracy. The other two are paired differences from the plain panel and from the earliest-arm rule, each with an interval at 90 percent from a bootstrap of 500 resamples of whole blocks. These intervals belong to no declared family and are not adjusted for the number of comparisons in the table, so they describe and carry no verdict. Table~\ref{tab:fmt-readers-methods-api} is the same table for run 12.}
\label{tab:fmt-readers-methods}
% generated by paper/build_tables.py from the report keys named above each row (readers); it reads paper/data/run11-simulated-readers
\begin{tabular}{p{4.6cm}rrr}
\toprule
model and method & share matched & minus the plain panel & minus the earliest-arm rule \\
\midrule
\multicolumn{4}{@{}l}{\emph{Opus}} \\
% src: paper/data/run11-simulated-readers/simulated-readers-claude-opus-5-5-20261106-2dc461f52ea4.json plain.accuracy; paper/data/run11-simulated-readers/simulated-readers-claude-opus-5-5-20261106-2dc461f52ea4.json readings.plain_minus_always_reference.estimate; paper/data/run11-simulated-readers/simulated-readers-claude-opus-5-5-20261106-2dc461f52ea4.json readings.plain_minus_always_reference.lo; paper/data/run11-simulated-readers/simulated-readers-claude-opus-5-5-20261106-2dc461f52ea4.json readings.plain_minus_always_reference.hi
\quad plain panel & 0.5742 &  & 0.0041 (-0.0117, 0.0217) \\
% src: paper/data/run11-simulated-readers/simulated-readers-claude-opus-5-5-20261106-2dc461f52ea4.json methods.direction_map.accuracy; paper/data/run11-simulated-readers/simulated-readers-claude-opus-5-5-20261106-2dc461f52ea4.json methods.direction_map.minus_plain.estimate; paper/data/run11-simulated-readers/simulated-readers-claude-opus-5-5-20261106-2dc461f52ea4.json methods.direction_map.minus_plain.lo; paper/data/run11-simulated-readers/simulated-readers-claude-opus-5-5-20261106-2dc461f52ea4.json methods.direction_map.minus_plain.hi; paper/data/run11-simulated-readers/simulated-readers-claude-opus-5-5-20261106-2dc461f52ea4.json methods.direction_map.minus_always_reference.estimate; paper/data/run11-simulated-readers/simulated-readers-claude-opus-5-5-20261106-2dc461f52ea4.json methods.direction_map.minus_always_reference.lo; paper/data/run11-simulated-readers/simulated-readers-claude-opus-5-5-20261106-2dc461f52ea4.json methods.direction_map.minus_always_reference.hi
\quad 1. direction map & 0.5840 & 0.0098 (-0.0013, 0.0200) & 0.0139 (0.0039, 0.0247) \\
% src: paper/data/run11-simulated-readers/simulated-readers-claude-opus-5-5-20261106-2dc461f52ea4.json methods.reader_weights.accuracy; paper/data/run11-simulated-readers/simulated-readers-claude-opus-5-5-20261106-2dc461f52ea4.json methods.reader_weights.minus_plain.estimate; paper/data/run11-simulated-readers/simulated-readers-claude-opus-5-5-20261106-2dc461f52ea4.json methods.reader_weights.minus_plain.lo; paper/data/run11-simulated-readers/simulated-readers-claude-opus-5-5-20261106-2dc461f52ea4.json methods.reader_weights.minus_plain.hi; paper/data/run11-simulated-readers/simulated-readers-claude-opus-5-5-20261106-2dc461f52ea4.json methods.reader_weights.minus_always_reference.estimate; paper/data/run11-simulated-readers/simulated-readers-claude-opus-5-5-20261106-2dc461f52ea4.json methods.reader_weights.minus_always_reference.lo; paper/data/run11-simulated-readers/simulated-readers-claude-opus-5-5-20261106-2dc461f52ea4.json methods.reader_weights.minus_always_reference.hi
\quad 2. reader weights & 0.5882 & 0.0140 (0.0041, 0.0241) & 0.0181 (0.0002, 0.0368) \\
% src: paper/data/run11-simulated-readers/simulated-readers-claude-opus-5-5-20261106-2dc461f52ea4.json methods.audience_matching.accuracy; paper/data/run11-simulated-readers/simulated-readers-claude-opus-5-5-20261106-2dc461f52ea4.json methods.audience_matching.minus_plain.estimate; paper/data/run11-simulated-readers/simulated-readers-claude-opus-5-5-20261106-2dc461f52ea4.json methods.audience_matching.minus_plain.lo; paper/data/run11-simulated-readers/simulated-readers-claude-opus-5-5-20261106-2dc461f52ea4.json methods.audience_matching.minus_plain.hi; paper/data/run11-simulated-readers/simulated-readers-claude-opus-5-5-20261106-2dc461f52ea4.json methods.audience_matching.minus_always_reference.estimate; paper/data/run11-simulated-readers/simulated-readers-claude-opus-5-5-20261106-2dc461f52ea4.json methods.audience_matching.minus_always_reference.lo; paper/data/run11-simulated-readers/simulated-readers-claude-opus-5-5-20261106-2dc461f52ea4.json methods.audience_matching.minus_always_reference.hi
\quad 3. audience matching & 0.5716 & -0.0026 (-0.0139, 0.0066) & 0.0015 (-0.0165, 0.0218) \\
% src: paper/data/run11-simulated-readers/simulated-readers-claude-opus-5-5-20261106-2dc461f52ea4.json methods.error_correction.accuracy; paper/data/run11-simulated-readers/simulated-readers-claude-opus-5-5-20261106-2dc461f52ea4.json methods.error_correction.minus_plain.estimate; paper/data/run11-simulated-readers/simulated-readers-claude-opus-5-5-20261106-2dc461f52ea4.json methods.error_correction.minus_plain.lo; paper/data/run11-simulated-readers/simulated-readers-claude-opus-5-5-20261106-2dc461f52ea4.json methods.error_correction.minus_plain.hi; paper/data/run11-simulated-readers/simulated-readers-claude-opus-5-5-20261106-2dc461f52ea4.json methods.error_correction.minus_always_reference.estimate; paper/data/run11-simulated-readers/simulated-readers-claude-opus-5-5-20261106-2dc461f52ea4.json methods.error_correction.minus_always_reference.lo; paper/data/run11-simulated-readers/simulated-readers-claude-opus-5-5-20261106-2dc461f52ea4.json methods.error_correction.minus_always_reference.hi
\quad 4. error correction & 0.5819 & 0.0077 (-0.0016, 0.0157) & 0.0118 (-0.0062, 0.0312) \\
% src: paper/data/run11-simulated-readers/simulated-readers-claude-opus-5-5-20261106-2dc461f52ea4.json methods.text_with_panel.accuracy; paper/data/run11-simulated-readers/simulated-readers-claude-opus-5-5-20261106-2dc461f52ea4.json methods.text_with_panel.minus_plain.estimate; paper/data/run11-simulated-readers/simulated-readers-claude-opus-5-5-20261106-2dc461f52ea4.json methods.text_with_panel.minus_plain.lo; paper/data/run11-simulated-readers/simulated-readers-claude-opus-5-5-20261106-2dc461f52ea4.json methods.text_with_panel.minus_plain.hi; paper/data/run11-simulated-readers/simulated-readers-claude-opus-5-5-20261106-2dc461f52ea4.json methods.text_with_panel.minus_always_reference.estimate; paper/data/run11-simulated-readers/simulated-readers-claude-opus-5-5-20261106-2dc461f52ea4.json methods.text_with_panel.minus_always_reference.lo; paper/data/run11-simulated-readers/simulated-readers-claude-opus-5-5-20261106-2dc461f52ea4.json methods.text_with_panel.minus_always_reference.hi
\quad 5. text classifier with the panel & 0.5687 & -0.0054 (-0.0195, 0.0077) & -0.0013 (-0.0153, 0.0122) \\
% src: paper/data/run11-simulated-readers/simulated-readers-claude-opus-5-5-20261106-2dc461f52ea4.json methods.examples_in_prompt.accuracy; paper/data/run11-simulated-readers/simulated-readers-claude-opus-5-5-20261106-2dc461f52ea4.json methods.examples_in_prompt.minus_plain.estimate; paper/data/run11-simulated-readers/simulated-readers-claude-opus-5-5-20261106-2dc461f52ea4.json methods.examples_in_prompt.minus_plain.lo; paper/data/run11-simulated-readers/simulated-readers-claude-opus-5-5-20261106-2dc461f52ea4.json methods.examples_in_prompt.minus_plain.hi; paper/data/run11-simulated-readers/simulated-readers-claude-opus-5-5-20261106-2dc461f52ea4.json methods.examples_in_prompt.minus_always_reference.estimate; paper/data/run11-simulated-readers/simulated-readers-claude-opus-5-5-20261106-2dc461f52ea4.json methods.examples_in_prompt.minus_always_reference.lo; paper/data/run11-simulated-readers/simulated-readers-claude-opus-5-5-20261106-2dc461f52ea4.json methods.examples_in_prompt.minus_always_reference.hi
\quad 6. earlier results in the prompt & 0.5853 & 0.0111 (0.0019, 0.0217) & 0.0153 (-0.0022, 0.0345) \\
% src: paper/data/run11-simulated-readers/simulated-readers-claude-opus-5-5-20261106-2dc461f52ea4.json methods.audience_brief.accuracy; paper/data/run11-simulated-readers/simulated-readers-claude-opus-5-5-20261106-2dc461f52ea4.json methods.audience_brief.minus_plain.estimate; paper/data/run11-simulated-readers/simulated-readers-claude-opus-5-5-20261106-2dc461f52ea4.json methods.audience_brief.minus_plain.lo; paper/data/run11-simulated-readers/simulated-readers-claude-opus-5-5-20261106-2dc461f52ea4.json methods.audience_brief.minus_plain.hi; paper/data/run11-simulated-readers/simulated-readers-claude-opus-5-5-20261106-2dc461f52ea4.json methods.audience_brief.minus_always_reference.estimate; paper/data/run11-simulated-readers/simulated-readers-claude-opus-5-5-20261106-2dc461f52ea4.json methods.audience_brief.minus_always_reference.lo; paper/data/run11-simulated-readers/simulated-readers-claude-opus-5-5-20261106-2dc461f52ea4.json methods.audience_brief.minus_always_reference.hi
\quad 7. audience brief & 0.6122 & 0.0380 (0.0274, 0.0480) & 0.0422 (0.0260, 0.0575) \\
\midrule
\multicolumn{4}{@{}l}{\emph{Sonnet}} \\
% src: paper/data/run11-simulated-readers/simulated-readers-claude-sonnet-5-5-20261106-2dc461f52ea4.json plain.accuracy; paper/data/run11-simulated-readers/simulated-readers-claude-sonnet-5-5-20261106-2dc461f52ea4.json readings.plain_minus_always_reference.estimate; paper/data/run11-simulated-readers/simulated-readers-claude-sonnet-5-5-20261106-2dc461f52ea4.json readings.plain_minus_always_reference.lo; paper/data/run11-simulated-readers/simulated-readers-claude-sonnet-5-5-20261106-2dc461f52ea4.json readings.plain_minus_always_reference.hi
\quad plain panel & 0.5234 &  & -0.0467 (-0.0648, -0.0275) \\
% src: paper/data/run11-simulated-readers/simulated-readers-claude-sonnet-5-5-20261106-2dc461f52ea4.json methods.direction_map.accuracy; paper/data/run11-simulated-readers/simulated-readers-claude-sonnet-5-5-20261106-2dc461f52ea4.json methods.direction_map.minus_plain.estimate; paper/data/run11-simulated-readers/simulated-readers-claude-sonnet-5-5-20261106-2dc461f52ea4.json methods.direction_map.minus_plain.lo; paper/data/run11-simulated-readers/simulated-readers-claude-sonnet-5-5-20261106-2dc461f52ea4.json methods.direction_map.minus_plain.hi; paper/data/run11-simulated-readers/simulated-readers-claude-sonnet-5-5-20261106-2dc461f52ea4.json methods.direction_map.minus_always_reference.estimate; paper/data/run11-simulated-readers/simulated-readers-claude-sonnet-5-5-20261106-2dc461f52ea4.json methods.direction_map.minus_always_reference.lo; paper/data/run11-simulated-readers/simulated-readers-claude-sonnet-5-5-20261106-2dc461f52ea4.json methods.direction_map.minus_always_reference.hi
\quad 1. direction map & 0.5622 & 0.0388 (0.0233, 0.0550) & -0.0079 (-0.0143, -0.0016) \\
% src: paper/data/run11-simulated-readers/simulated-readers-claude-sonnet-5-5-20261106-2dc461f52ea4.json methods.reader_weights.accuracy; paper/data/run11-simulated-readers/simulated-readers-claude-sonnet-5-5-20261106-2dc461f52ea4.json methods.reader_weights.minus_plain.estimate; paper/data/run11-simulated-readers/simulated-readers-claude-sonnet-5-5-20261106-2dc461f52ea4.json methods.reader_weights.minus_plain.lo; paper/data/run11-simulated-readers/simulated-readers-claude-sonnet-5-5-20261106-2dc461f52ea4.json methods.reader_weights.minus_plain.hi; paper/data/run11-simulated-readers/simulated-readers-claude-sonnet-5-5-20261106-2dc461f52ea4.json methods.reader_weights.minus_always_reference.estimate; paper/data/run11-simulated-readers/simulated-readers-claude-sonnet-5-5-20261106-2dc461f52ea4.json methods.reader_weights.minus_always_reference.lo; paper/data/run11-simulated-readers/simulated-readers-claude-sonnet-5-5-20261106-2dc461f52ea4.json methods.reader_weights.minus_always_reference.hi
\quad 2. reader weights & 0.5598 & 0.0365 (0.0253, 0.0474) & -0.0102 (-0.0308, 0.0082) \\
% src: paper/data/run11-simulated-readers/simulated-readers-claude-sonnet-5-5-20261106-2dc461f52ea4.json methods.audience_matching.accuracy; paper/data/run11-simulated-readers/simulated-readers-claude-sonnet-5-5-20261106-2dc461f52ea4.json methods.audience_matching.minus_plain.estimate; paper/data/run11-simulated-readers/simulated-readers-claude-sonnet-5-5-20261106-2dc461f52ea4.json methods.audience_matching.minus_plain.lo; paper/data/run11-simulated-readers/simulated-readers-claude-sonnet-5-5-20261106-2dc461f52ea4.json methods.audience_matching.minus_plain.hi; paper/data/run11-simulated-readers/simulated-readers-claude-sonnet-5-5-20261106-2dc461f52ea4.json methods.audience_matching.minus_always_reference.estimate; paper/data/run11-simulated-readers/simulated-readers-claude-sonnet-5-5-20261106-2dc461f52ea4.json methods.audience_matching.minus_always_reference.lo; paper/data/run11-simulated-readers/simulated-readers-claude-sonnet-5-5-20261106-2dc461f52ea4.json methods.audience_matching.minus_always_reference.hi
\quad 3. audience matching & 0.5357 & 0.0123 (0.0031, 0.0211) & -0.0344 (-0.0515, -0.0189) \\
% src: paper/data/run11-simulated-readers/simulated-readers-claude-sonnet-5-5-20261106-2dc461f52ea4.json methods.error_correction.accuracy; paper/data/run11-simulated-readers/simulated-readers-claude-sonnet-5-5-20261106-2dc461f52ea4.json methods.error_correction.minus_plain.estimate; paper/data/run11-simulated-readers/simulated-readers-claude-sonnet-5-5-20261106-2dc461f52ea4.json methods.error_correction.minus_plain.lo; paper/data/run11-simulated-readers/simulated-readers-claude-sonnet-5-5-20261106-2dc461f52ea4.json methods.error_correction.minus_plain.hi; paper/data/run11-simulated-readers/simulated-readers-claude-sonnet-5-5-20261106-2dc461f52ea4.json methods.error_correction.minus_always_reference.estimate; paper/data/run11-simulated-readers/simulated-readers-claude-sonnet-5-5-20261106-2dc461f52ea4.json methods.error_correction.minus_always_reference.lo; paper/data/run11-simulated-readers/simulated-readers-claude-sonnet-5-5-20261106-2dc461f52ea4.json methods.error_correction.minus_always_reference.hi
\quad 4. error correction & 0.5484 & 0.0251 (0.0169, 0.0329) & -0.0216 (-0.0404, -0.0042) \\
% src: paper/data/run11-simulated-readers/simulated-readers-claude-sonnet-5-5-20261106-2dc461f52ea4.json methods.text_with_panel.accuracy; paper/data/run11-simulated-readers/simulated-readers-claude-sonnet-5-5-20261106-2dc461f52ea4.json methods.text_with_panel.minus_plain.estimate; paper/data/run11-simulated-readers/simulated-readers-claude-sonnet-5-5-20261106-2dc461f52ea4.json methods.text_with_panel.minus_plain.lo; paper/data/run11-simulated-readers/simulated-readers-claude-sonnet-5-5-20261106-2dc461f52ea4.json methods.text_with_panel.minus_plain.hi; paper/data/run11-simulated-readers/simulated-readers-claude-sonnet-5-5-20261106-2dc461f52ea4.json methods.text_with_panel.minus_always_reference.estimate; paper/data/run11-simulated-readers/simulated-readers-claude-sonnet-5-5-20261106-2dc461f52ea4.json methods.text_with_panel.minus_always_reference.lo; paper/data/run11-simulated-readers/simulated-readers-claude-sonnet-5-5-20261106-2dc461f52ea4.json methods.text_with_panel.minus_always_reference.hi
\quad 5. text classifier with the panel & 0.5484 & 0.0251 (0.0098, 0.0402) & -0.0216 (-0.0367, -0.0060) \\
% src: paper/data/run11-simulated-readers/simulated-readers-claude-sonnet-5-5-20261106-2dc461f52ea4.json methods.examples_in_prompt.accuracy; paper/data/run11-simulated-readers/simulated-readers-claude-sonnet-5-5-20261106-2dc461f52ea4.json methods.examples_in_prompt.minus_plain.estimate; paper/data/run11-simulated-readers/simulated-readers-claude-sonnet-5-5-20261106-2dc461f52ea4.json methods.examples_in_prompt.minus_plain.lo; paper/data/run11-simulated-readers/simulated-readers-claude-sonnet-5-5-20261106-2dc461f52ea4.json methods.examples_in_prompt.minus_plain.hi; paper/data/run11-simulated-readers/simulated-readers-claude-sonnet-5-5-20261106-2dc461f52ea4.json methods.examples_in_prompt.minus_always_reference.estimate; paper/data/run11-simulated-readers/simulated-readers-claude-sonnet-5-5-20261106-2dc461f52ea4.json methods.examples_in_prompt.minus_always_reference.lo; paper/data/run11-simulated-readers/simulated-readers-claude-sonnet-5-5-20261106-2dc461f52ea4.json methods.examples_in_prompt.minus_always_reference.hi
\quad 6. earlier results in the prompt & 0.5419 & 0.0185 (0.0073, 0.0301) & -0.0282 (-0.0487, -0.0102) \\
% src: paper/data/run11-simulated-readers/simulated-readers-claude-sonnet-5-5-20261106-2dc461f52ea4.json methods.audience_brief.accuracy; paper/data/run11-simulated-readers/simulated-readers-claude-sonnet-5-5-20261106-2dc461f52ea4.json methods.audience_brief.minus_plain.estimate; paper/data/run11-simulated-readers/simulated-readers-claude-sonnet-5-5-20261106-2dc461f52ea4.json methods.audience_brief.minus_plain.lo; paper/data/run11-simulated-readers/simulated-readers-claude-sonnet-5-5-20261106-2dc461f52ea4.json methods.audience_brief.minus_plain.hi; paper/data/run11-simulated-readers/simulated-readers-claude-sonnet-5-5-20261106-2dc461f52ea4.json methods.audience_brief.minus_always_reference.estimate; paper/data/run11-simulated-readers/simulated-readers-claude-sonnet-5-5-20261106-2dc461f52ea4.json methods.audience_brief.minus_always_reference.lo; paper/data/run11-simulated-readers/simulated-readers-claude-sonnet-5-5-20261106-2dc461f52ea4.json methods.audience_brief.minus_always_reference.hi
\quad 7. audience brief & 0.5828 & 0.0594 (0.0455, 0.0715) & 0.0127 (-0.0067, 0.0308) \\
\midrule
\multicolumn{4}{@{}l}{\emph{Haiku}} \\
% src: paper/data/run11-simulated-readers/simulated-readers-claude-haiku-5-5-20261106-2dc461f52ea4.json plain.accuracy; paper/data/run11-simulated-readers/simulated-readers-claude-haiku-5-5-20261106-2dc461f52ea4.json readings.plain_minus_always_reference.estimate; paper/data/run11-simulated-readers/simulated-readers-claude-haiku-5-5-20261106-2dc461f52ea4.json readings.plain_minus_always_reference.lo; paper/data/run11-simulated-readers/simulated-readers-claude-haiku-5-5-20261106-2dc461f52ea4.json readings.plain_minus_always_reference.hi
\quad plain panel & 0.5144 &  & -0.0556 (-0.0747, -0.0385) \\
% src: paper/data/run11-simulated-readers/simulated-readers-claude-haiku-5-5-20261106-2dc461f52ea4.json methods.direction_map.accuracy; paper/data/run11-simulated-readers/simulated-readers-claude-haiku-5-5-20261106-2dc461f52ea4.json methods.direction_map.minus_plain.estimate; paper/data/run11-simulated-readers/simulated-readers-claude-haiku-5-5-20261106-2dc461f52ea4.json methods.direction_map.minus_plain.lo; paper/data/run11-simulated-readers/simulated-readers-claude-haiku-5-5-20261106-2dc461f52ea4.json methods.direction_map.minus_plain.hi; paper/data/run11-simulated-readers/simulated-readers-claude-haiku-5-5-20261106-2dc461f52ea4.json methods.direction_map.minus_always_reference.estimate; paper/data/run11-simulated-readers/simulated-readers-claude-haiku-5-5-20261106-2dc461f52ea4.json methods.direction_map.minus_always_reference.lo; paper/data/run11-simulated-readers/simulated-readers-claude-haiku-5-5-20261106-2dc461f52ea4.json methods.direction_map.minus_always_reference.hi
\quad 1. direction map & 0.5669 & 0.0525 (0.0357, 0.0685) & -0.0032 (-0.0117, 0.0042) \\
% src: paper/data/run11-simulated-readers/simulated-readers-claude-haiku-5-5-20261106-2dc461f52ea4.json methods.reader_weights.accuracy; paper/data/run11-simulated-readers/simulated-readers-claude-haiku-5-5-20261106-2dc461f52ea4.json methods.reader_weights.minus_plain.estimate; paper/data/run11-simulated-readers/simulated-readers-claude-haiku-5-5-20261106-2dc461f52ea4.json methods.reader_weights.minus_plain.lo; paper/data/run11-simulated-readers/simulated-readers-claude-haiku-5-5-20261106-2dc461f52ea4.json methods.reader_weights.minus_plain.hi; paper/data/run11-simulated-readers/simulated-readers-claude-haiku-5-5-20261106-2dc461f52ea4.json methods.reader_weights.minus_always_reference.estimate; paper/data/run11-simulated-readers/simulated-readers-claude-haiku-5-5-20261106-2dc461f52ea4.json methods.reader_weights.minus_always_reference.lo; paper/data/run11-simulated-readers/simulated-readers-claude-haiku-5-5-20261106-2dc461f52ea4.json methods.reader_weights.minus_always_reference.hi
\quad 2. reader weights & 0.5375 & 0.0231 (0.0101, 0.0349) & -0.0325 (-0.0531, -0.0145) \\
% src: paper/data/run11-simulated-readers/simulated-readers-claude-haiku-5-5-20261106-2dc461f52ea4.json methods.audience_matching.accuracy; paper/data/run11-simulated-readers/simulated-readers-claude-haiku-5-5-20261106-2dc461f52ea4.json methods.audience_matching.minus_plain.estimate; paper/data/run11-simulated-readers/simulated-readers-claude-haiku-5-5-20261106-2dc461f52ea4.json methods.audience_matching.minus_plain.lo; paper/data/run11-simulated-readers/simulated-readers-claude-haiku-5-5-20261106-2dc461f52ea4.json methods.audience_matching.minus_plain.hi; paper/data/run11-simulated-readers/simulated-readers-claude-haiku-5-5-20261106-2dc461f52ea4.json methods.audience_matching.minus_always_reference.estimate; paper/data/run11-simulated-readers/simulated-readers-claude-haiku-5-5-20261106-2dc461f52ea4.json methods.audience_matching.minus_always_reference.lo; paper/data/run11-simulated-readers/simulated-readers-claude-haiku-5-5-20261106-2dc461f52ea4.json methods.audience_matching.minus_always_reference.hi
\quad 3. audience matching & 0.5020 & -0.0124 (-0.0231, -0.0002) & -0.0681 (-0.0865, -0.0489) \\
% src: paper/data/run11-simulated-readers/simulated-readers-claude-haiku-5-5-20261106-2dc461f52ea4.json methods.error_correction.accuracy; paper/data/run11-simulated-readers/simulated-readers-claude-haiku-5-5-20261106-2dc461f52ea4.json methods.error_correction.minus_plain.estimate; paper/data/run11-simulated-readers/simulated-readers-claude-haiku-5-5-20261106-2dc461f52ea4.json methods.error_correction.minus_plain.lo; paper/data/run11-simulated-readers/simulated-readers-claude-haiku-5-5-20261106-2dc461f52ea4.json methods.error_correction.minus_plain.hi; paper/data/run11-simulated-readers/simulated-readers-claude-haiku-5-5-20261106-2dc461f52ea4.json methods.error_correction.minus_always_reference.estimate; paper/data/run11-simulated-readers/simulated-readers-claude-haiku-5-5-20261106-2dc461f52ea4.json methods.error_correction.minus_always_reference.lo; paper/data/run11-simulated-readers/simulated-readers-claude-haiku-5-5-20261106-2dc461f52ea4.json methods.error_correction.minus_always_reference.hi
\quad 4. error correction & 0.5360 & 0.0216 (0.0129, 0.0298) & -0.0340 (-0.0525, -0.0175) \\
% src: paper/data/run11-simulated-readers/simulated-readers-claude-haiku-5-5-20261106-2dc461f52ea4.json methods.text_with_panel.accuracy; paper/data/run11-simulated-readers/simulated-readers-claude-haiku-5-5-20261106-2dc461f52ea4.json methods.text_with_panel.minus_plain.estimate; paper/data/run11-simulated-readers/simulated-readers-claude-haiku-5-5-20261106-2dc461f52ea4.json methods.text_with_panel.minus_plain.lo; paper/data/run11-simulated-readers/simulated-readers-claude-haiku-5-5-20261106-2dc461f52ea4.json methods.text_with_panel.minus_plain.hi; paper/data/run11-simulated-readers/simulated-readers-claude-haiku-5-5-20261106-2dc461f52ea4.json methods.text_with_panel.minus_always_reference.estimate; paper/data/run11-simulated-readers/simulated-readers-claude-haiku-5-5-20261106-2dc461f52ea4.json methods.text_with_panel.minus_always_reference.lo; paper/data/run11-simulated-readers/simulated-readers-claude-haiku-5-5-20261106-2dc461f52ea4.json methods.text_with_panel.minus_always_reference.hi
\quad 5. text classifier with the panel & 0.5499 & 0.0355 (0.0216, 0.0518) & -0.0201 (-0.0355, -0.0038) \\
% src: paper/data/run11-simulated-readers/simulated-readers-claude-haiku-5-5-20261106-2dc461f52ea4.json methods.examples_in_prompt.accuracy; paper/data/run11-simulated-readers/simulated-readers-claude-haiku-5-5-20261106-2dc461f52ea4.json methods.examples_in_prompt.minus_plain.estimate; paper/data/run11-simulated-readers/simulated-readers-claude-haiku-5-5-20261106-2dc461f52ea4.json methods.examples_in_prompt.minus_plain.lo; paper/data/run11-simulated-readers/simulated-readers-claude-haiku-5-5-20261106-2dc461f52ea4.json methods.examples_in_prompt.minus_plain.hi; paper/data/run11-simulated-readers/simulated-readers-claude-haiku-5-5-20261106-2dc461f52ea4.json methods.examples_in_prompt.minus_always_reference.estimate; paper/data/run11-simulated-readers/simulated-readers-claude-haiku-5-5-20261106-2dc461f52ea4.json methods.examples_in_prompt.minus_always_reference.lo; paper/data/run11-simulated-readers/simulated-readers-claude-haiku-5-5-20261106-2dc461f52ea4.json methods.examples_in_prompt.minus_always_reference.hi
\quad 6. earlier results in the prompt & 0.5189 & 0.0045 (-0.0095, 0.0180) & -0.0511 (-0.0703, -0.0348) \\
% src: paper/data/run11-simulated-readers/simulated-readers-claude-haiku-5-5-20261106-2dc461f52ea4.json methods.audience_brief.accuracy; paper/data/run11-simulated-readers/simulated-readers-claude-haiku-5-5-20261106-2dc461f52ea4.json methods.audience_brief.minus_plain.estimate; paper/data/run11-simulated-readers/simulated-readers-claude-haiku-5-5-20261106-2dc461f52ea4.json methods.audience_brief.minus_plain.lo; paper/data/run11-simulated-readers/simulated-readers-claude-haiku-5-5-20261106-2dc461f52ea4.json methods.audience_brief.minus_plain.hi; paper/data/run11-simulated-readers/simulated-readers-claude-haiku-5-5-20261106-2dc461f52ea4.json methods.audience_brief.minus_always_reference.estimate; paper/data/run11-simulated-readers/simulated-readers-claude-haiku-5-5-20261106-2dc461f52ea4.json methods.audience_brief.minus_always_reference.lo; paper/data/run11-simulated-readers/simulated-readers-claude-haiku-5-5-20261106-2dc461f52ea4.json methods.audience_brief.minus_always_reference.hi
\quad 7. audience brief & 0.5671 & 0.0527 (0.0392, 0.0687) & -0.0029 (-0.0206, 0.0147) \\
\bottomrule
\end{tabular}

\end{table*}

\begin{table}[t]\centering\footnotesize
% src: tables/fmt-readers-transport.tex; paper/data/run12-simulated-readers-api/simulated-readers-transport-claude-opus-5-5-20261106-2dc461f52ea4.json plain.accuracy_api_minus_tool.resamples 500
\caption{The plain panels of the two runs, which send the same requests by two routes. The difference in direction accuracy has a 90 percent bootstrap interval from 500 resamples of blocks and belongs to no family. The share with the same direction is taken over the contrasts on which neither panel is tied.}
\label{tab:fmt-readers-transport}
\resizebox{\columnwidth}{!}{% generated by paper/build_tables.py from the report keys named above each row (readers); it reads paper/data/run12-simulated-readers-api
\begin{tabular}{lrrrrr}
\toprule
model & run 11 & run 12 & run 12 minus run 11 & same direction, & contrasts \\
 & & & & share & compared \\
\midrule
% src: paper/data/run12-simulated-readers-api/simulated-readers-transport-claude-opus-5-5-20261106-2dc461f52ea4.json plain.accuracy_tool; paper/data/run12-simulated-readers-api/simulated-readers-transport-claude-opus-5-5-20261106-2dc461f52ea4.json plain.accuracy_api; paper/data/run12-simulated-readers-api/simulated-readers-transport-claude-opus-5-5-20261106-2dc461f52ea4.json plain.accuracy_api_minus_tool.estimate; paper/data/run12-simulated-readers-api/simulated-readers-transport-claude-opus-5-5-20261106-2dc461f52ea4.json plain.accuracy_api_minus_tool.lo; paper/data/run12-simulated-readers-api/simulated-readers-transport-claude-opus-5-5-20261106-2dc461f52ea4.json plain.accuracy_api_minus_tool.hi; paper/data/run12-simulated-readers-api/simulated-readers-transport-claude-opus-5-5-20261106-2dc461f52ea4.json plain.agree_share; paper/data/run12-simulated-readers-api/simulated-readers-transport-claude-opus-5-5-20261106-2dc461f52ea4.json plain.n
Opus & 0.5742 & 0.5742 & 0.0000 (-0.0085, 0.0081) & 0.852 & 5,962 \\
% src: paper/data/run12-simulated-readers-api/simulated-readers-transport-claude-sonnet-5-5-20261106-2dc461f52ea4.json plain.accuracy_tool; paper/data/run12-simulated-readers-api/simulated-readers-transport-claude-sonnet-5-5-20261106-2dc461f52ea4.json plain.accuracy_api; paper/data/run12-simulated-readers-api/simulated-readers-transport-claude-sonnet-5-5-20261106-2dc461f52ea4.json plain.accuracy_api_minus_tool.estimate; paper/data/run12-simulated-readers-api/simulated-readers-transport-claude-sonnet-5-5-20261106-2dc461f52ea4.json plain.accuracy_api_minus_tool.lo; paper/data/run12-simulated-readers-api/simulated-readers-transport-claude-sonnet-5-5-20261106-2dc461f52ea4.json plain.accuracy_api_minus_tool.hi; paper/data/run12-simulated-readers-api/simulated-readers-transport-claude-sonnet-5-5-20261106-2dc461f52ea4.json plain.agree_share; paper/data/run12-simulated-readers-api/simulated-readers-transport-claude-sonnet-5-5-20261106-2dc461f52ea4.json plain.n
Sonnet & 0.5234 & 0.5279 & 0.0045 (-0.0040, 0.0134) & 0.812 & 5,964 \\
% src: paper/data/run12-simulated-readers-api/simulated-readers-transport-claude-haiku-5-5-20261106-2dc461f52ea4.json plain.accuracy_tool; paper/data/run12-simulated-readers-api/simulated-readers-transport-claude-haiku-5-5-20261106-2dc461f52ea4.json plain.accuracy_api; paper/data/run12-simulated-readers-api/simulated-readers-transport-claude-haiku-5-5-20261106-2dc461f52ea4.json plain.accuracy_api_minus_tool.estimate; paper/data/run12-simulated-readers-api/simulated-readers-transport-claude-haiku-5-5-20261106-2dc461f52ea4.json plain.accuracy_api_minus_tool.lo; paper/data/run12-simulated-readers-api/simulated-readers-transport-claude-haiku-5-5-20261106-2dc461f52ea4.json plain.accuracy_api_minus_tool.hi; paper/data/run12-simulated-readers-api/simulated-readers-transport-claude-haiku-5-5-20261106-2dc461f52ea4.json plain.agree_share; paper/data/run12-simulated-readers-api/simulated-readers-transport-claude-haiku-5-5-20261106-2dc461f52ea4.json plain.n
Haiku & 0.5144 & 0.5223 & 0.0079 (-0.0035, 0.0206) & 0.671 & 5,964 \\
\bottomrule
\end{tabular}
}
\end{table}

% src: tables/fmt-readers-methods.tex Opus row 7. audience brief 0.6122, minus the earliest-arm rule 0.0422 (0.0260, 0.0575), Sonnet 0.0127 (-0.0067, 0.0308), Haiku -0.0029 (-0.0206, 0.0147); tables/fmt-readers-methods-api.tex Opus row 7. audience brief 0.6170, minus the earliest-arm rule 0.0469 (0.0300, 0.0623), Haiku 0.5401, minus the earliest-arm rule -0.0300 (-0.0472, -0.0131); tables/fmt-readers-readings.tex the direction ceiling 0.8301; tables/fmt-readers-methods.tex row 3. audience matching, minus the plain panel: Haiku -0.0124 (-0.0231, -0.0002); tables/fmt-readers-methods-api.tex row 3. audience matching, minus the plain panel: Opus -0.0091 (-0.0166, -0.0010), Haiku -0.0141 (-0.0254, -0.0045); run 11 and run 12 are the names of the two run folders
\paragraph{The seven methods.} Table~\ref{tab:fmt-readers-methods} sets the methods side by side. One stands out for the primary model: with the audience brief its panel is right on 0.612 of the contrasts in run 11 and on 0.617 in run 12. Its difference from the earliest-arm rule is 0.042 in run 11, with an interval that lies above zero, and run 12 agrees. We report this as a description and not as a finding. The comparison was not one of the three readings and was chosen after the table was seen. Its interval lies above zero in both runs and is not adjusted for the seven methods and three models of the table. No adjustment would remove the fact that we picked the best of seven methods after seeing them. The model writes a new note in each run, so for this method run 12 is a second draw and not a repeat of the same requests. The design expected this method to be among the weakest, so what the table shows was not predicted. For the smallest model the same method is below the earliest-arm rule in run 12. Even the best figure here is far from the ceiling of 0.830. Audience matching, the reweighting that uses no outcome of a test, has the lowest or the second lowest accuracy of the seven methods for every model in both runs. For the smallest model in run 12 it is below the plain panel, with an unadjusted interval that lies clear of zero. For the same model in run 11 and for the primary model in run 12 the estimate is also negative, but the upper end of each interval is within the margin of Section~\ref{sec:protocol} of zero, so we do not read those two. An alignment may therefore lower accuracy as well as raise it. These rows describe and carry no verdict, and we apply the margin to them so that no row is read more freely than the paper's other tables.

% src: tables/fmt-readers-transport.tex Opus 0.5742, 0.5742, 0.0000 (-0.0085, 0.0081), same direction 0.852; Sonnet 0.0045 (-0.0040, 0.0134); Haiku 0.0079 (-0.0035, 0.0206), same direction 0.671; tables/fmt-readers-more.tex row the two halves of the panel give the same direction, share: Opus 0.750, row contrasts on which the plain panel is tied: Opus 3; tables/fmt-readers-more-api.tex row contrasts on which the plain panel is tied: Opus 1 (the replies of the two runs differ); paper/data/run12-simulated-readers-api/simulated-readers-transport-claude-opus-5-5-20261106-2dc461f52ea4.json plain.accuracy_api_minus_tool.estimate 0.0; src/mootcourt/simulators/readers/protocol.py HALVES (two panels of five readers); paper/data/run12-simulated-readers-api/README.md, The transport (thinking and effort left at each model's default); no request was sent twice by one route; run 11 is the name of a run folder
\paragraph{The two runs against each other.} For the primary model the plain panel's accuracy is exactly the same in the two runs. That is a coincidence: the replies differ, as the rest of this paragraph shows. The interval of the difference runs from minus 0.0085 to 0.0081 (Table~\ref{tab:fmt-readers-transport}). By the rule fixed beforehand, the context that the command line tool adds did not move the plain panel's accuracy by an amount this design can detect. The two routes differ in that context and possibly in default settings as well (the paragraph on how the calls were made), so this sentence is about the change of route as a whole. The corresponding interval contains zero for each of the two smaller models. Equal accuracy is not equal answers: the two panels of the primary model give the same direction on 0.852 of the contrasts, and those of the smallest model on 0.671. Within run 11, two halves of the primary model's panel agree on 0.750 of the contrasts. These two kinds of figure are not like for like, because the second compares panels of five readers. We sent no request twice by the same route, so we cannot say how much of the disagreement between the runs is chance. A panel of ten readers is thus a noisy instrument, and a difference between two runs that lies inside the intervals of Table~\ref{tab:fmt-readers-transport} should not be read.

% src: tables/fmt-readers-probes.tex (completion: winner minus other, each within two standard errors of zero in both runs; familiarity: the headline rated more familiar won, share, 0.509 to 0.534, above one half in all six cases; row familiarity: that share minus one half (90 percent interval), lower end above zero by more than 0.0010 for Sonnet in both runs, 0.0122 and 0.0026, and for Opus in run 12, 0.0026, and below zero in the other three cases; the margin 0.0010 is that of Section sec:protocol; familiarity: same direction as the plain panel, share; named publisher: 0.557 to 0.595)
\paragraph{Earlier exposure to the archive.} Three probes ask whether a model recalls the winning headline of a test better than a losing one, which it could if the publisher ran the winner afterwards (Table~\ref{tab:fmt-readers-probes}). Given the first four words of a headline, a model completes the winning headline of a test no better than another headline of the same test. This first probe gives no sign that winners are recalled better. A model's ratings of how familiar a headline is point to the winner on 0.509 to 0.534 of the contrasts probed. That share is above one half for every model in both runs. An interval over resampled blocks, added after the results, lies above one half in three of the six cases: for the middle model in both runs and for the primary model in run 12. In the other three cases it contains one half. The intervals are not adjusted for the six cases. By the rule fixed with these intervals we read the ratings as pointing to the winner more often than chance for the middle model. For the primary model one run of two supports that, which is the pattern we treat as not established for the second reading. The open model of run 13 received the same three probes; their results are in the reports of that run and are not printed. A model may also rate an appealing or a typical headline as familiar, so this second probe cannot separate recall from judgement. Asked directly which headline of a named publisher's pair won, each model is right more often than a coin would be, which it can do by judging the headlines and which is therefore no evidence either way. One probe thus gives no sign of recall, one points to the winner more often than chance for one of the three hosted models in both runs and for another in one, and the third cannot tell. None of them can rule recall out, and every statement of this section is made with contamination not excluded.

% src: src/mootcourt/simulators/readers/protocol.py SYSTEM_PROMPT, READER_TEMPLATE (copied word for word; the numbers 2013, 2015, 0 and 100 are part of the copy); no measured number
\paragraph{The instructions and the prompt.} Every reader call carries the same instructions:
\begin{quote}\small
You are playing one person who is browsing the web at some time between 2013 and 2015. You will be given a short description of that person and then a feed of story headlines from a news and entertainment site, numbered, one per line. For each headline, give the probability from 0 to 100 that this person would click on it at this moment. Judge each headline as this person would on a quick scroll. Most people click on few of the stories they scroll past. Reply with one JSON object and nothing else.
\end{quote}
The prompt then gives the reader's profile, the numbered headlines of the feed, and one line that asks for a list of as many numbers as there are headlines. A reply that does not hold exactly that list is asked for again, up to three attempts in all, and is then recorded as missing. A reader with a missing feed in a block is dropped from that block, so that every arm of a test is averaged over the same readers.

% src: src/mootcourt/simulators/readers/protocol.py PROFILES (the ten profiles, copied word for word; the ages 18, 19, 24, 25, 34, 35, 44, 45, 54, 55 and 64 and the five minutes are part of the copy); no measured number
\paragraph{The ten profiles.} The profiles were written once, before any call, with the paper's other results known.
\begin{enumerate}\small
\item You are between 19 and 24 and a college student. You like music, campus politics and funny videos. You are putting off an assignment and are a little bored. You are on your phone, and a friend shared this in your social feed.
\item You are between 25 and 34 and work in an office. You follow national news, technology and personal finance. You are on a short lunch break and want something quick. You are at your desktop, and the story came up in your social feed.
\item You are between 25 and 34 and have two small children. You care about parenting, health and stories that end well. It is late in the evening and you are tired. You are on your phone, and a relative shared this in your social feed.
\item You are between 35 and 44 and teach at a public school. You care about education, science and your local community. You are relaxed after dinner. You are at your desktop, and the story arrived in an email newsletter from the site.
\item You are between 35 and 44 and work as a nurse. You like stories about health, animals and ordinary people. You have five minutes of a break and are in a hurry. You are on your phone, and the story came up in your social feed.
\item You are between 45 and 54 and manage a small business. You follow the economy, politics and sports. You are sceptical of anything that sounds exaggerated. You are at your desktop, and someone you follow posted the link.
\item You are between 45 and 54 and at home with two teenagers. You care about family, faith and cooking. You are unhurried and in a good mood. You are at your desktop, and a neighbour shared this in your social feed.
\item You are between 55 and 64 and recently retired. You like gardening, history and news about your grandchildren's generation. You have plenty of time. You are at your desktop, and a friend forwarded the story by email.
\item You are between 18 and 24 and work shifts in a shop. You like celebrities, video games and comedy. You are restless and scrolling fast. You are on your phone, and the story came up in your social feed.
\item You are between 25 and 34 and work for a charity. You care about the environment, equal rights and documentaries. You are curious and willing to read something serious. You are on your phone, and you are on the site's front page.
\end{enumerate}

% src: src/mootcourt/simulators/readers/protocol.py EXAMPLES_HEADER, BRIEF_TEMPLATE, BRIEF_PREFACE, N_EXAMPLES 10, N_BRIEF_UNITS 200, NOTE_WORDS 150; src/mootcourt/simulators/readers/incontext.py example_units (drawn once per origin: every feed of the origin shows the same ten), brief_units (every visible unit when fewer than 200 exist), clip_note (the first 150 words of a longer note); paper/data/run11-simulated-readers/plan.json counts.by_origin.2014-01-01.earlier_units_visible 163 (fewer than 200 at the second scored origin; in words only), counts.calls_by_kind.brief_note 6 (one draw and one note at each of the six origins after the first); tables/fmt-readers-more.tex and tables/fmt-readers-more-api.tex (calls planned 9,007); paper/data/run11-simulated-readers/simulated-readers-claude-opus-5-5-20261106-2dc461f52ea4.json counts.calls_by_kind.plain 3080, counts.calls_by_kind.examples 2860, counts.calls_by_kind.brief 2860, counts.calls_by_kind.brief_note 6, counts.calls_by_kind.probe_completion 60, counts.calls_by_kind.probe_familiarity 42, counts.calls_by_kind.probe_pairs 99 (201 probe calls: 60 + 42 + 99), counts.calls_planned 9007
\paragraph{The two methods that change the prompt.} For method 6 the prompt opens with ten earlier tests, drawn with a fixed seed from the tests visible at the origin, each as its headlines with the share of visitors who clicked. For method 7 the model first receives up to 200 visible earlier tests in the same form, all of them when fewer are visible, as at the second scored origin. It is asked for a note of at most 150 words on what this audience clicked and passed over, without quoting a headline. A longer note is cut at 150 words. Each reader's prompt at that origin then opens with the note. One draw serves every feed and reader of an origin, so each of the two methods rests on six draws, one for each origin after the first. Both methods reuse the blocks and feeds of the plain panel. One model's calls number 9,007. The plain panel takes 3,080 of them and each of the two methods 2,860. The notes take 6 and the probes 201.

% src: src/mootcourt/simulators/readers/align.py direction_map (LogisticRegression of the variant winning on the panel difference), _fit_reader_weights (softmax weights; binomial log likelihood of clicks given impressions, mean per arm; log odds an intercept plus a positive slope times the weighted mean of the centred scores; WEIGHT_PENALTY 1.0 times the squared distance from equal weights), _fit_audience_matching (squared gap between the weighted mean raw probability and the pooled click rate; MATCH_PENALTY 0.01 times the same distance), error_correction (a logistic curve from panel score to click rate; least squares of its error, weighted by impressions, on a constant and four traits), traits (the count of words, a question mark, a digit, the word you), _reader_matrix (the arms that all ten readers scored), visible_before; src/mootcourt/simulators/readers/textfeature.py (one more column; a row with no value gets zero; fitted on the visible rows of every cell); paper/data/run11-simulated-readers/README.md, Deviations from the plan (the objective of the reader weights); the numbers 1 to 5 are the methods' numbers; no measured number
\paragraph{How methods 1 to 5 are fitted.} Each fit uses the earlier scored tests that Section~\ref{sec:readers} names, and only the arms that all ten readers scored. Method 1 is a logistic regression, with the default penalty of its library, of whether the variant won on the panel's difference between variant and reference. Its direction is the side of one half on which its probability falls. Method 2 models the clicks of each earlier arm as binomial given the arm's impressions. The log odds of a click are an intercept plus a positive slope times the weighted mean of the ten readers' centred probabilities, and the weights are positive and sum to one. The fit maximizes the mean log likelihood per arm, less a penalty of strength 1 on the squared distance of the weights from equal weights. This is the objective as changed before the first call (Appendix~\ref{sec:analysis-history}). Method 3 takes each reader's mean uncentred probability over the earlier arms. Its weights, again positive with sum one, minimize the squared gap between the weighted mean of those levels and the pooled click rate of the earlier arms, with the same penalty at strength 0.01. Method 4 first fits a logistic curve from an arm's panel score to its click rate. It then regresses the curve's error, by least squares weighted by impressions, on a constant and four traits of the headline: its count of words, whether it holds a question mark, whether it holds a digit, and whether it holds the word ``you''. A scored arm receives the curve's value plus that correction, and the direction is the sign of the difference between two arms. Method 5 adds the panel's difference as one more column to the text classifier of Section~\ref{sec:method}, which is fitted on the visible rows of every cell. Training rows without a panel answer take the value zero in that column.

% src: paper/data/run11-simulated-readers/README.md (the tool at one version; its flags; an empty directory; no tool and no server; a subscription login, under which the tool puts its own context around the instructions: the model's name and knowledge cutoff and the present date among it; no record of tokens; requests that took a second or third attempt 15, 20 and 59, in words only; the full suite ran once before the review fixes and only parts were rerun afterwards); docs/simulated-readers.md (the tool's context cannot be switched off under a subscription login); paper/data/run12-simulated-readers-api/README.md (four request fields; the defaults left in place; the vendor documents the default effort as medium for claude-opus-5-5 and claude-haiku-5-5 and as high for claude-sonnet-5-5; requests that took a second or third attempt 23, 18 and 100, in words only; the token counts and the cost of 132.13 dollars at list prices; four invalid replies of the smallest model, all of method 7); tables/fmt-readers-more-api.tex rows usable replies and invalid after three attempts; run 11 and run 12 are the names of the two run folders
\paragraph{How the calls were made.} In run 11 each call started the vendor's command line tool in an empty directory with its tools switched off, no settings read and our instructions in place of its own. The tool ran under a subscription login, and under such a login it still adds a short context of its own that cannot be switched off. That context includes the present date, as the paragraph on the two runs describes, and the model's name and knowledge cutoff. The run stored no token counts. In run 12 each call is one request to the vendor's programming interface that holds the model's name, our instructions, the prompt and a limit on the length of the reply. No thinking or effort setting was sent in run 11 or run 12, so each route used its own default, and we did not check that the two defaults are the same. For the programming interface the vendor documents a default effort of medium for the largest and the smallest model and of high for the middle one. A difference between models is therefore a difference between models at their defaults. Every request and reply is stored, and the scoring reads the stored replies and calls no model. In run 11 every planned call returned a usable reply. In both runs some requests needed a second or third attempt before the reply had the required form, most often for the smallest model, and the stored reply is the first one of that form. The READMEs of the run folders give the counts. In run 12 four calls of the smallest model, all of method 7, returned no usable reply after three attempts, and the rule above dropped those readers from their blocks (Table~\ref{tab:fmt-readers-more-api}). The second run cost about 132 dollars at the vendor's list prices.

% src: paper/data/run13-readers-open-model/README.md (run 13, the name of its run folder; Qwen/Qwen3.8-27B at a recorded revision, full precision, one card of 80 GB; the plan folder of run 11 unchanged; no sampling setting; the cost rule of 40 dollars, the projection of 65 to 118 dollars at the default, so medium, set on the server and not in each request; paper/data/run13-readers-open-model/cost-projection.txt (the projection comes from sixteen calls on an invented feed of 25 lines and not from the canary of twenty calls that the plan names; both are in the run folder); docs/superpowers/plans/2026-11-21-readers-quarters-and-open-model.md, Part B, Requests (a canary of twenty reader calls on invented headlines); one draw; about 22 dollars; four calls made at a second start of the stage; 8,769 headlines embedded by Qwen/Qwen3-Embedding-8B at dimension 4,096, no instruction text, unit length); tables/fmt-readers-open-model.tex rows planned calls 9,007 and 9,007, calls of the direct judgment 616 and 616; src/mootcourt/simulators/readers/direct.py DIRECT_TEMPLATE, DIRECT_BATCH 10; src/mootcourt/simulators/readers/embedfeature.py embedding_directions (C of 1.0; the design: the difference of the two embeddings, the variant's embedding, one column per cell, and for the version with the panel the panel difference divided by ten); docs/superpowers/plans/2026-11-21-readers-quarters-and-open-model.md, Part B
\paragraph{The run with an open model (run 13).} The plan of this run was written after the results of runs 11 and 12 and before any reply of the open model or any embedding existed. The model is Qwen3.8-27B, served at full precision on one rented graphics card, and the run folder records its revision. It received the requests of run 11 unchanged, 9,007 calls, with no sampling setting sent. One rule on thinking was fixed beforehand: the model's default applies unless a trial of twenty calls on invented headlines projects the run above 40 US dollars, and then a medium setting applies to every request. The run departs from its plan in two ways here. The projection that decided was taken from a second trial on a longer invented feed, which is closer to the length of a real feed, and not from the trial the plan names; the run folder holds both, and the projection from the named trial alone would have been lower. The medium setting was set once on the server and not in each request, because the server does not read it from a request. Both trials read invented headlines only, so the choice could not depend on a result. Every planned call returned a usable reply. Four calls returned none at the first start of the stage and were made when the stage was started again. The model was run once, so its replies are one draw, and the run cost about 22 dollars.

For the direct judgment a call lists ten contrasts of one block. Each is shown as two headlines, A and B, in a drawn order, under the sentence that the two were written for the same story, and the model names for each pair the headline that more people would click. The 616 calls cover every contrast once, and the answer is right when it names the headline with the higher click rate.

The embedding classifier mirrors the text classifier of the main run. Each headline is embedded once by an open embedding model, Qwen3-Embedding-8B, with no instruction text, and the vector is scaled to unit length. At each origin a logistic regression with the library's default regularization is fitted on the admitted rows visible at that origin. Its columns are the difference of the two headlines' embeddings, the variant's embedding and one column per cell. The version with a panel adds one column, the panel's difference divided by ten, as method 5 does. The embedding model was also trained after the archive was public, so its vectors may carry exposure to these headlines. Nothing in the classifier was tuned, so a null from it would not show that no text model can beat the earliest-arm rule, and its lead on resampled blocks is one model's.

\begin{table*}[t]\centering\footnotesize
% src: tables/fmt-readers-quarters.tex; paper/data/run11-simulated-readers/simulated-readers-quarters-claude-opus-5-5-20261106-2dc461f52ea4.json readings.weights_minus_plain.quarters.origins 7, readings.weights_minus_plain.quarters.family_size 3, level 0.9 (the unadjusted level, which the table does not print); docs/superpowers/plans/2026-11-21-readers-quarters-and-open-model.md (the estimator, fixed after the results of both runs)
\caption{The three readings of the simulated readers with the seven quarters as the units, added after the results of both runs. Each cell is the mean of the seven quarterly differences in direction accuracy with a Student-t interval at the Bonferroni level for the three readings, and the column \emph{above} counts the quarters whose difference is above zero. For the second and third readings the first quarter is zero by construction, because no earlier scored test exists there, and it still counts as a unit. The quarters are of unequal size and each counts once. The interval treats the seven values as exchangeable, with the limits that Section~\ref{sec:ladder-results} states. No cell carries a verdict word: the fixed words of Table~\ref{tab:fmt-readers-readings} rest on intervals over resampled blocks.}
\label{tab:fmt-readers-quarters}
% generated by paper/build_tables.py from the report keys named above each row (readers); it reads paper/data/run11-simulated-readers and paper/data/run12-simulated-readers-api
\begin{tabular}{p{2.2cm}rrrr}
\toprule
reading and model & run 11 & above & run 12 & above \\
\midrule
\multicolumn{5}{@{}l}{\emph{1. Plain panel minus the earliest-arm rule}} \\
% src: paper/data/run11-simulated-readers/simulated-readers-quarters-claude-opus-5-5-20261106-2dc461f52ea4.json readings.plain_minus_always_reference.quarters.estimate; paper/data/run11-simulated-readers/simulated-readers-quarters-claude-opus-5-5-20261106-2dc461f52ea4.json readings.plain_minus_always_reference.quarters.adjusted_lo; paper/data/run11-simulated-readers/simulated-readers-quarters-claude-opus-5-5-20261106-2dc461f52ea4.json readings.plain_minus_always_reference.quarters.adjusted_hi; paper/data/run11-simulated-readers/simulated-readers-quarters-claude-opus-5-5-20261106-2dc461f52ea4.json readings.plain_minus_always_reference.quarters.origins_above_zero; paper/data/run12-simulated-readers-api/simulated-readers-quarters-claude-opus-5-5-run12-api-20261106-2dc461f52ea4.json readings.plain_minus_always_reference.quarters.estimate; paper/data/run12-simulated-readers-api/simulated-readers-quarters-claude-opus-5-5-run12-api-20261106-2dc461f52ea4.json readings.plain_minus_always_reference.quarters.adjusted_lo; paper/data/run12-simulated-readers-api/simulated-readers-quarters-claude-opus-5-5-run12-api-20261106-2dc461f52ea4.json readings.plain_minus_always_reference.quarters.adjusted_hi; paper/data/run12-simulated-readers-api/simulated-readers-quarters-claude-opus-5-5-run12-api-20261106-2dc461f52ea4.json readings.plain_minus_always_reference.quarters.origins_above_zero
\quad Opus & 0.0070 (-0.0165, 0.0306) & 4 & 0.0104 (-0.0152, 0.0361) & 4 \\
% src: paper/data/run11-simulated-readers/simulated-readers-quarters-claude-sonnet-5-5-20261106-2dc461f52ea4.json readings.plain_minus_always_reference.quarters.estimate; paper/data/run11-simulated-readers/simulated-readers-quarters-claude-sonnet-5-5-20261106-2dc461f52ea4.json readings.plain_minus_always_reference.quarters.adjusted_lo; paper/data/run11-simulated-readers/simulated-readers-quarters-claude-sonnet-5-5-20261106-2dc461f52ea4.json readings.plain_minus_always_reference.quarters.adjusted_hi; paper/data/run11-simulated-readers/simulated-readers-quarters-claude-sonnet-5-5-20261106-2dc461f52ea4.json readings.plain_minus_always_reference.quarters.origins_above_zero; paper/data/run12-simulated-readers-api/simulated-readers-quarters-claude-sonnet-5-5-run12-api-20261106-2dc461f52ea4.json readings.plain_minus_always_reference.quarters.estimate; paper/data/run12-simulated-readers-api/simulated-readers-quarters-claude-sonnet-5-5-run12-api-20261106-2dc461f52ea4.json readings.plain_minus_always_reference.quarters.adjusted_lo; paper/data/run12-simulated-readers-api/simulated-readers-quarters-claude-sonnet-5-5-run12-api-20261106-2dc461f52ea4.json readings.plain_minus_always_reference.quarters.adjusted_hi; paper/data/run12-simulated-readers-api/simulated-readers-quarters-claude-sonnet-5-5-run12-api-20261106-2dc461f52ea4.json readings.plain_minus_always_reference.quarters.origins_above_zero
\quad Sonnet & -0.0470 (-0.0859, -0.0082) & 0 & -0.0401 (-0.0635, -0.0166) & 0 \\
% src: paper/data/run11-simulated-readers/simulated-readers-quarters-claude-haiku-5-5-20261106-2dc461f52ea4.json readings.plain_minus_always_reference.quarters.estimate; paper/data/run11-simulated-readers/simulated-readers-quarters-claude-haiku-5-5-20261106-2dc461f52ea4.json readings.plain_minus_always_reference.quarters.adjusted_lo; paper/data/run11-simulated-readers/simulated-readers-quarters-claude-haiku-5-5-20261106-2dc461f52ea4.json readings.plain_minus_always_reference.quarters.adjusted_hi; paper/data/run11-simulated-readers/simulated-readers-quarters-claude-haiku-5-5-20261106-2dc461f52ea4.json readings.plain_minus_always_reference.quarters.origins_above_zero; paper/data/run12-simulated-readers-api/simulated-readers-quarters-claude-haiku-5-5-run12-api-20261106-2dc461f52ea4.json readings.plain_minus_always_reference.quarters.estimate; paper/data/run12-simulated-readers-api/simulated-readers-quarters-claude-haiku-5-5-run12-api-20261106-2dc461f52ea4.json readings.plain_minus_always_reference.quarters.adjusted_lo; paper/data/run12-simulated-readers-api/simulated-readers-quarters-claude-haiku-5-5-run12-api-20261106-2dc461f52ea4.json readings.plain_minus_always_reference.quarters.adjusted_hi; paper/data/run12-simulated-readers-api/simulated-readers-quarters-claude-haiku-5-5-run12-api-20261106-2dc461f52ea4.json readings.plain_minus_always_reference.quarters.origins_above_zero
\quad Haiku & -0.0581 (-0.1003, -0.0158) & 0 & -0.0534 (-0.1010, -0.0058) & 0 \\
\midrule
\multicolumn{5}{@{}l}{\emph{2. Panel with reader weights minus the plain panel}} \\
% src: paper/data/run11-simulated-readers/simulated-readers-quarters-claude-opus-5-5-20261106-2dc461f52ea4.json readings.weights_minus_plain.quarters.estimate; paper/data/run11-simulated-readers/simulated-readers-quarters-claude-opus-5-5-20261106-2dc461f52ea4.json readings.weights_minus_plain.quarters.adjusted_lo; paper/data/run11-simulated-readers/simulated-readers-quarters-claude-opus-5-5-20261106-2dc461f52ea4.json readings.weights_minus_plain.quarters.adjusted_hi; paper/data/run11-simulated-readers/simulated-readers-quarters-claude-opus-5-5-20261106-2dc461f52ea4.json readings.weights_minus_plain.quarters.origins_above_zero; paper/data/run12-simulated-readers-api/simulated-readers-quarters-claude-opus-5-5-run12-api-20261106-2dc461f52ea4.json readings.weights_minus_plain.quarters.estimate; paper/data/run12-simulated-readers-api/simulated-readers-quarters-claude-opus-5-5-run12-api-20261106-2dc461f52ea4.json readings.weights_minus_plain.quarters.adjusted_lo; paper/data/run12-simulated-readers-api/simulated-readers-quarters-claude-opus-5-5-run12-api-20261106-2dc461f52ea4.json readings.weights_minus_plain.quarters.adjusted_hi; paper/data/run12-simulated-readers-api/simulated-readers-quarters-claude-opus-5-5-run12-api-20261106-2dc461f52ea4.json readings.weights_minus_plain.quarters.origins_above_zero
\quad Opus & 0.0123 (-0.0035, 0.0281) & 5 & 0.0080 (-0.0195, 0.0355) & 4 \\
% src: paper/data/run11-simulated-readers/simulated-readers-quarters-claude-sonnet-5-5-20261106-2dc461f52ea4.json readings.weights_minus_plain.quarters.estimate; paper/data/run11-simulated-readers/simulated-readers-quarters-claude-sonnet-5-5-20261106-2dc461f52ea4.json readings.weights_minus_plain.quarters.adjusted_lo; paper/data/run11-simulated-readers/simulated-readers-quarters-claude-sonnet-5-5-20261106-2dc461f52ea4.json readings.weights_minus_plain.quarters.adjusted_hi; paper/data/run11-simulated-readers/simulated-readers-quarters-claude-sonnet-5-5-20261106-2dc461f52ea4.json readings.weights_minus_plain.quarters.origins_above_zero; paper/data/run12-simulated-readers-api/simulated-readers-quarters-claude-sonnet-5-5-run12-api-20261106-2dc461f52ea4.json readings.weights_minus_plain.quarters.estimate; paper/data/run12-simulated-readers-api/simulated-readers-quarters-claude-sonnet-5-5-run12-api-20261106-2dc461f52ea4.json readings.weights_minus_plain.quarters.adjusted_lo; paper/data/run12-simulated-readers-api/simulated-readers-quarters-claude-sonnet-5-5-run12-api-20261106-2dc461f52ea4.json readings.weights_minus_plain.quarters.adjusted_hi; paper/data/run12-simulated-readers-api/simulated-readers-quarters-claude-sonnet-5-5-run12-api-20261106-2dc461f52ea4.json readings.weights_minus_plain.quarters.origins_above_zero
\quad Sonnet & 0.0350 (0.0036, 0.0663) & 6 & 0.0265 (0.0006, 0.0524) & 5 \\
% src: paper/data/run11-simulated-readers/simulated-readers-quarters-claude-haiku-5-5-20261106-2dc461f52ea4.json readings.weights_minus_plain.quarters.estimate; paper/data/run11-simulated-readers/simulated-readers-quarters-claude-haiku-5-5-20261106-2dc461f52ea4.json readings.weights_minus_plain.quarters.adjusted_lo; paper/data/run11-simulated-readers/simulated-readers-quarters-claude-haiku-5-5-20261106-2dc461f52ea4.json readings.weights_minus_plain.quarters.adjusted_hi; paper/data/run11-simulated-readers/simulated-readers-quarters-claude-haiku-5-5-20261106-2dc461f52ea4.json readings.weights_minus_plain.quarters.origins_above_zero; paper/data/run12-simulated-readers-api/simulated-readers-quarters-claude-haiku-5-5-run12-api-20261106-2dc461f52ea4.json readings.weights_minus_plain.quarters.estimate; paper/data/run12-simulated-readers-api/simulated-readers-quarters-claude-haiku-5-5-run12-api-20261106-2dc461f52ea4.json readings.weights_minus_plain.quarters.adjusted_lo; paper/data/run12-simulated-readers-api/simulated-readers-quarters-claude-haiku-5-5-run12-api-20261106-2dc461f52ea4.json readings.weights_minus_plain.quarters.adjusted_hi; paper/data/run12-simulated-readers-api/simulated-readers-quarters-claude-haiku-5-5-run12-api-20261106-2dc461f52ea4.json readings.weights_minus_plain.quarters.origins_above_zero
\quad Haiku & 0.0254 (0.0079, 0.0430) & 6 & 0.0039 (-0.0161, 0.0238) & 3 \\
\midrule
\multicolumn{5}{@{}l}{\emph{3. Text classifier with the panel minus the text classifier alone}} \\
% src: paper/data/run11-simulated-readers/simulated-readers-quarters-claude-opus-5-5-20261106-2dc461f52ea4.json readings.text_with_panel_minus_text.quarters.estimate; paper/data/run11-simulated-readers/simulated-readers-quarters-claude-opus-5-5-20261106-2dc461f52ea4.json readings.text_with_panel_minus_text.quarters.adjusted_lo; paper/data/run11-simulated-readers/simulated-readers-quarters-claude-opus-5-5-20261106-2dc461f52ea4.json readings.text_with_panel_minus_text.quarters.adjusted_hi; paper/data/run11-simulated-readers/simulated-readers-quarters-claude-opus-5-5-20261106-2dc461f52ea4.json readings.text_with_panel_minus_text.quarters.origins_above_zero; paper/data/run12-simulated-readers-api/simulated-readers-quarters-claude-opus-5-5-run12-api-20261106-2dc461f52ea4.json readings.text_with_panel_minus_text.quarters.estimate; paper/data/run12-simulated-readers-api/simulated-readers-quarters-claude-opus-5-5-run12-api-20261106-2dc461f52ea4.json readings.text_with_panel_minus_text.quarters.adjusted_lo; paper/data/run12-simulated-readers-api/simulated-readers-quarters-claude-opus-5-5-run12-api-20261106-2dc461f52ea4.json readings.text_with_panel_minus_text.quarters.adjusted_hi; paper/data/run12-simulated-readers-api/simulated-readers-quarters-claude-opus-5-5-run12-api-20261106-2dc461f52ea4.json readings.text_with_panel_minus_text.quarters.origins_above_zero
\quad Opus & 0.0246 (0.0013, 0.0479) & 6 & 0.0225 (0.0026, 0.0424) & 6 \\
% src: paper/data/run11-simulated-readers/simulated-readers-quarters-claude-sonnet-5-5-20261106-2dc461f52ea4.json readings.text_with_panel_minus_text.quarters.estimate; paper/data/run11-simulated-readers/simulated-readers-quarters-claude-sonnet-5-5-20261106-2dc461f52ea4.json readings.text_with_panel_minus_text.quarters.adjusted_lo; paper/data/run11-simulated-readers/simulated-readers-quarters-claude-sonnet-5-5-20261106-2dc461f52ea4.json readings.text_with_panel_minus_text.quarters.adjusted_hi; paper/data/run11-simulated-readers/simulated-readers-quarters-claude-sonnet-5-5-20261106-2dc461f52ea4.json readings.text_with_panel_minus_text.quarters.origins_above_zero; paper/data/run12-simulated-readers-api/simulated-readers-quarters-claude-sonnet-5-5-run12-api-20261106-2dc461f52ea4.json readings.text_with_panel_minus_text.quarters.estimate; paper/data/run12-simulated-readers-api/simulated-readers-quarters-claude-sonnet-5-5-run12-api-20261106-2dc461f52ea4.json readings.text_with_panel_minus_text.quarters.adjusted_lo; paper/data/run12-simulated-readers-api/simulated-readers-quarters-claude-sonnet-5-5-run12-api-20261106-2dc461f52ea4.json readings.text_with_panel_minus_text.quarters.adjusted_hi; paper/data/run12-simulated-readers-api/simulated-readers-quarters-claude-sonnet-5-5-run12-api-20261106-2dc461f52ea4.json readings.text_with_panel_minus_text.quarters.origins_above_zero
\quad Sonnet & 0.0020 (-0.0083, 0.0124) & 2 & 0.0021 (-0.0101, 0.0143) & 3 \\
% src: paper/data/run11-simulated-readers/simulated-readers-quarters-claude-haiku-5-5-20261106-2dc461f52ea4.json readings.text_with_panel_minus_text.quarters.estimate; paper/data/run11-simulated-readers/simulated-readers-quarters-claude-haiku-5-5-20261106-2dc461f52ea4.json readings.text_with_panel_minus_text.quarters.adjusted_lo; paper/data/run11-simulated-readers/simulated-readers-quarters-claude-haiku-5-5-20261106-2dc461f52ea4.json readings.text_with_panel_minus_text.quarters.adjusted_hi; paper/data/run11-simulated-readers/simulated-readers-quarters-claude-haiku-5-5-20261106-2dc461f52ea4.json readings.text_with_panel_minus_text.quarters.origins_above_zero; paper/data/run12-simulated-readers-api/simulated-readers-quarters-claude-haiku-5-5-run12-api-20261106-2dc461f52ea4.json readings.text_with_panel_minus_text.quarters.estimate; paper/data/run12-simulated-readers-api/simulated-readers-quarters-claude-haiku-5-5-run12-api-20261106-2dc461f52ea4.json readings.text_with_panel_minus_text.quarters.adjusted_lo; paper/data/run12-simulated-readers-api/simulated-readers-quarters-claude-haiku-5-5-run12-api-20261106-2dc461f52ea4.json readings.text_with_panel_minus_text.quarters.adjusted_hi; paper/data/run12-simulated-readers-api/simulated-readers-quarters-claude-haiku-5-5-run12-api-20261106-2dc461f52ea4.json readings.text_with_panel_minus_text.quarters.origins_above_zero
\quad Haiku & 0.0043 (-0.0038, 0.0123) & 3 & 0.0088 (-0.0050, 0.0226) & 4 \\
\bottomrule
\end{tabular}

\end{table*}

\begin{table*}[t]\centering\footnotesize
% src: tables/fmt-readers-open-model.tex; paper/data/run13-readers-open-model/simulated-readers-open-model-qwen3.8-27b-run13-open-20261106-2dc461f52ea4.json level 0.9, direct.n 5966, direct.minus_plain.blocks.resamples 500, direct.minus_plain.quarters.origins 7, note (added after the results of runs 11 and 12; no entry carries a verdict word); paper/data/run13-readers-open-model/simulated-readers-qwen3.8-27b-run13-open-20261106-2dc461f52ea4.json counts.blocks 63, primary false; run 13, 11 and 12 are the names of the run folders
\caption{Run 13, added after the results of runs 11 and 12: an open model's plain panel and its direct judgment, and a classifier on embeddings of the headlines, on the 5,966 contrasts of the headline-only cell. The upper rows give direction accuracy. Each difference is given twice: with an interval at 90 percent from 500 resamples of whole blocks, and as the mean of the seven quarterly differences with a Student-t interval at 90 percent. No interval is adjusted for the number of comparisons, and no row carries a verdict word. The last rows count the calls.}
\label{tab:fmt-readers-open-model}
% generated by paper/build_tables.py from the report keys named above each row (readers); it reads paper/data/run13-readers-open-model
\begin{tabular}{p{6.6cm}rr}
\toprule
measure & blocks resampled & quarters as units \\
\midrule
\multicolumn{3}{@{}l}{\emph{Share of contrasts whose direction is matched}} \\
% src: paper/data/run13-readers-open-model/simulated-readers-qwen3.8-27b-run13-open-20261106-2dc461f52ea4.json baselines.always_reference.accuracy
\quad the earliest-arm rule & 0.5701 &  \\
% src: paper/data/run13-readers-open-model/simulated-readers-qwen3.8-27b-run13-open-20261106-2dc461f52ea4.json baselines.text_lr.accuracy
\quad the text classifier of the main run & 0.5484 &  \\
% src: paper/data/run13-readers-open-model/simulated-readers-open-model-qwen3.8-27b-run13-open-20261106-2dc461f52ea4.json embedding.accuracy
\quad the embedding classifier & 0.6012 &  \\
% src: paper/data/run13-readers-open-model/simulated-readers-qwen3.8-27b-run13-open-20261106-2dc461f52ea4.json plain.accuracy
\quad the open model's plain panel & 0.5208 &  \\
% src: paper/data/run13-readers-open-model/simulated-readers-open-model-qwen3.8-27b-run13-open-20261106-2dc461f52ea4.json direct.accuracy
\quad the open model's direct judgment & 0.5607 &  \\
\midrule
\multicolumn{3}{@{}l}{\emph{The open model: difference, with blocks resampled and with the quarters as the units}} \\
% src: paper/data/run13-readers-open-model/simulated-readers-open-model-qwen3.8-27b-run13-open-20261106-2dc461f52ea4.json direct.minus_plain.blocks.estimate; paper/data/run13-readers-open-model/simulated-readers-open-model-qwen3.8-27b-run13-open-20261106-2dc461f52ea4.json direct.minus_plain.blocks.lo; paper/data/run13-readers-open-model/simulated-readers-open-model-qwen3.8-27b-run13-open-20261106-2dc461f52ea4.json direct.minus_plain.blocks.hi; paper/data/run13-readers-open-model/simulated-readers-open-model-qwen3.8-27b-run13-open-20261106-2dc461f52ea4.json direct.minus_plain.quarters.estimate; paper/data/run13-readers-open-model/simulated-readers-open-model-qwen3.8-27b-run13-open-20261106-2dc461f52ea4.json direct.minus_plain.quarters.lo; paper/data/run13-readers-open-model/simulated-readers-open-model-qwen3.8-27b-run13-open-20261106-2dc461f52ea4.json direct.minus_plain.quarters.hi
\quad direct judgment minus its plain panel & 0.0399 (0.0242, 0.0558) & 0.0253 (-0.0224, 0.0731) \\
% src: paper/data/run13-readers-open-model/simulated-readers-open-model-qwen3.8-27b-run13-open-20261106-2dc461f52ea4.json direct.minus_always_reference.blocks.estimate; paper/data/run13-readers-open-model/simulated-readers-open-model-qwen3.8-27b-run13-open-20261106-2dc461f52ea4.json direct.minus_always_reference.blocks.lo; paper/data/run13-readers-open-model/simulated-readers-open-model-qwen3.8-27b-run13-open-20261106-2dc461f52ea4.json direct.minus_always_reference.blocks.hi; paper/data/run13-readers-open-model/simulated-readers-open-model-qwen3.8-27b-run13-open-20261106-2dc461f52ea4.json direct.minus_always_reference.quarters.estimate; paper/data/run13-readers-open-model/simulated-readers-open-model-qwen3.8-27b-run13-open-20261106-2dc461f52ea4.json direct.minus_always_reference.quarters.lo; paper/data/run13-readers-open-model/simulated-readers-open-model-qwen3.8-27b-run13-open-20261106-2dc461f52ea4.json direct.minus_always_reference.quarters.hi
\quad direct judgment minus the earliest-arm rule & -0.0094 (-0.0303, 0.0105) & -0.0201 (-0.0764, 0.0362) \\
% src: paper/data/run13-readers-open-model/simulated-readers-open-model-qwen3.8-27b-run13-open-20261106-2dc461f52ea4.json direct.plain_minus_always_reference.blocks.estimate; paper/data/run13-readers-open-model/simulated-readers-open-model-qwen3.8-27b-run13-open-20261106-2dc461f52ea4.json direct.plain_minus_always_reference.blocks.lo; paper/data/run13-readers-open-model/simulated-readers-open-model-qwen3.8-27b-run13-open-20261106-2dc461f52ea4.json direct.plain_minus_always_reference.blocks.hi; paper/data/run13-readers-open-model/simulated-readers-open-model-qwen3.8-27b-run13-open-20261106-2dc461f52ea4.json direct.plain_minus_always_reference.quarters.estimate; paper/data/run13-readers-open-model/simulated-readers-open-model-qwen3.8-27b-run13-open-20261106-2dc461f52ea4.json direct.plain_minus_always_reference.quarters.lo; paper/data/run13-readers-open-model/simulated-readers-open-model-qwen3.8-27b-run13-open-20261106-2dc461f52ea4.json direct.plain_minus_always_reference.quarters.hi
\quad plain panel minus the earliest-arm rule & -0.0493 (-0.0667, -0.0314) & -0.0454 (-0.0641, -0.0267) \\
% src: paper/data/run13-readers-open-model/simulated-readers-quarters-qwen3.8-27b-run13-open-20261106-2dc461f52ea4.json readings.weights_minus_plain.blocks.estimate; paper/data/run13-readers-open-model/simulated-readers-quarters-qwen3.8-27b-run13-open-20261106-2dc461f52ea4.json readings.weights_minus_plain.blocks.lo; paper/data/run13-readers-open-model/simulated-readers-quarters-qwen3.8-27b-run13-open-20261106-2dc461f52ea4.json readings.weights_minus_plain.blocks.hi; paper/data/run13-readers-open-model/simulated-readers-quarters-qwen3.8-27b-run13-open-20261106-2dc461f52ea4.json readings.weights_minus_plain.quarters.estimate; paper/data/run13-readers-open-model/simulated-readers-quarters-qwen3.8-27b-run13-open-20261106-2dc461f52ea4.json readings.weights_minus_plain.quarters.lo; paper/data/run13-readers-open-model/simulated-readers-quarters-qwen3.8-27b-run13-open-20261106-2dc461f52ea4.json readings.weights_minus_plain.quarters.hi
\quad reader weights minus the plain panel & 0.0213 (0.0097, 0.0333) & 0.0198 (0.0006, 0.0390) \\
% src: paper/data/run13-readers-open-model/simulated-readers-quarters-qwen3.8-27b-run13-open-20261106-2dc461f52ea4.json readings.text_with_panel_minus_text.blocks.estimate; paper/data/run13-readers-open-model/simulated-readers-quarters-qwen3.8-27b-run13-open-20261106-2dc461f52ea4.json readings.text_with_panel_minus_text.blocks.lo; paper/data/run13-readers-open-model/simulated-readers-quarters-qwen3.8-27b-run13-open-20261106-2dc461f52ea4.json readings.text_with_panel_minus_text.blocks.hi; paper/data/run13-readers-open-model/simulated-readers-quarters-qwen3.8-27b-run13-open-20261106-2dc461f52ea4.json readings.text_with_panel_minus_text.quarters.estimate; paper/data/run13-readers-open-model/simulated-readers-quarters-qwen3.8-27b-run13-open-20261106-2dc461f52ea4.json readings.text_with_panel_minus_text.quarters.lo; paper/data/run13-readers-open-model/simulated-readers-quarters-qwen3.8-27b-run13-open-20261106-2dc461f52ea4.json readings.text_with_panel_minus_text.quarters.hi
\quad text classifier with its panel minus without & 0.0022 (-0.0029, 0.0078) & 0.0051 (-0.0021, 0.0123) \\
\midrule
\multicolumn{3}{@{}l}{\emph{The embedding classifier: difference, read the same two ways}} \\
% src: paper/data/run13-readers-open-model/simulated-readers-open-model-qwen3.8-27b-run13-open-20261106-2dc461f52ea4.json embedding.minus_always_reference.blocks.estimate; paper/data/run13-readers-open-model/simulated-readers-open-model-qwen3.8-27b-run13-open-20261106-2dc461f52ea4.json embedding.minus_always_reference.blocks.lo; paper/data/run13-readers-open-model/simulated-readers-open-model-qwen3.8-27b-run13-open-20261106-2dc461f52ea4.json embedding.minus_always_reference.blocks.hi; paper/data/run13-readers-open-model/simulated-readers-open-model-qwen3.8-27b-run13-open-20261106-2dc461f52ea4.json embedding.minus_always_reference.quarters.estimate; paper/data/run13-readers-open-model/simulated-readers-open-model-qwen3.8-27b-run13-open-20261106-2dc461f52ea4.json embedding.minus_always_reference.quarters.lo; paper/data/run13-readers-open-model/simulated-readers-open-model-qwen3.8-27b-run13-open-20261106-2dc461f52ea4.json embedding.minus_always_reference.quarters.hi
\quad minus the earliest-arm rule & 0.0312 (0.0125, 0.0460) & 0.0197 (-0.0191, 0.0585) \\
% src: paper/data/run13-readers-open-model/simulated-readers-open-model-qwen3.8-27b-run13-open-20261106-2dc461f52ea4.json embedding.minus_text.blocks.estimate; paper/data/run13-readers-open-model/simulated-readers-open-model-qwen3.8-27b-run13-open-20261106-2dc461f52ea4.json embedding.minus_text.blocks.lo; paper/data/run13-readers-open-model/simulated-readers-open-model-qwen3.8-27b-run13-open-20261106-2dc461f52ea4.json embedding.minus_text.blocks.hi; paper/data/run13-readers-open-model/simulated-readers-open-model-qwen3.8-27b-run13-open-20261106-2dc461f52ea4.json embedding.minus_text.quarters.estimate; paper/data/run13-readers-open-model/simulated-readers-open-model-qwen3.8-27b-run13-open-20261106-2dc461f52ea4.json embedding.minus_text.quarters.lo; paper/data/run13-readers-open-model/simulated-readers-open-model-qwen3.8-27b-run13-open-20261106-2dc461f52ea4.json embedding.minus_text.quarters.hi
\quad minus the text classifier of the main run & 0.0528 (0.0368, 0.0677) & 0.0542 (0.0377, 0.0707) \\
% src: paper/data/run13-readers-open-model/simulated-readers-open-model-qwen3.8-27b-run13-open-20261106-2dc461f52ea4.json embedding.with_panel["opus-run11"].minus_without.blocks.estimate; paper/data/run13-readers-open-model/simulated-readers-open-model-qwen3.8-27b-run13-open-20261106-2dc461f52ea4.json embedding.with_panel["opus-run11"].minus_without.blocks.lo; paper/data/run13-readers-open-model/simulated-readers-open-model-qwen3.8-27b-run13-open-20261106-2dc461f52ea4.json embedding.with_panel["opus-run11"].minus_without.blocks.hi; paper/data/run13-readers-open-model/simulated-readers-open-model-qwen3.8-27b-run13-open-20261106-2dc461f52ea4.json embedding.with_panel["opus-run11"].minus_without.quarters.estimate; paper/data/run13-readers-open-model/simulated-readers-open-model-qwen3.8-27b-run13-open-20261106-2dc461f52ea4.json embedding.with_panel["opus-run11"].minus_without.quarters.lo; paper/data/run13-readers-open-model/simulated-readers-open-model-qwen3.8-27b-run13-open-20261106-2dc461f52ea4.json embedding.with_panel["opus-run11"].minus_without.quarters.hi
\quad with the primary model's panel of run 11, minus without & 0.0101 (0.0027, 0.0182) & 0.0154 (0.0007, 0.0302) \\
% src: paper/data/run13-readers-open-model/simulated-readers-open-model-qwen3.8-27b-run13-open-20261106-2dc461f52ea4.json embedding.with_panel["opus-run12"].minus_without.blocks.estimate; paper/data/run13-readers-open-model/simulated-readers-open-model-qwen3.8-27b-run13-open-20261106-2dc461f52ea4.json embedding.with_panel["opus-run12"].minus_without.blocks.lo; paper/data/run13-readers-open-model/simulated-readers-open-model-qwen3.8-27b-run13-open-20261106-2dc461f52ea4.json embedding.with_panel["opus-run12"].minus_without.blocks.hi; paper/data/run13-readers-open-model/simulated-readers-open-model-qwen3.8-27b-run13-open-20261106-2dc461f52ea4.json embedding.with_panel["opus-run12"].minus_without.quarters.estimate; paper/data/run13-readers-open-model/simulated-readers-open-model-qwen3.8-27b-run13-open-20261106-2dc461f52ea4.json embedding.with_panel["opus-run12"].minus_without.quarters.lo; paper/data/run13-readers-open-model/simulated-readers-open-model-qwen3.8-27b-run13-open-20261106-2dc461f52ea4.json embedding.with_panel["opus-run12"].minus_without.quarters.hi
\quad with the primary model's panel of run 12, minus without & 0.0112 (0.0042, 0.0192) & 0.0140 (0.0051, 0.0229) \\
% src: paper/data/run13-readers-open-model/simulated-readers-open-model-qwen3.8-27b-run13-open-20261106-2dc461f52ea4.json embedding.with_panel["open-run13"].minus_without.blocks.estimate; paper/data/run13-readers-open-model/simulated-readers-open-model-qwen3.8-27b-run13-open-20261106-2dc461f52ea4.json embedding.with_panel["open-run13"].minus_without.blocks.lo; paper/data/run13-readers-open-model/simulated-readers-open-model-qwen3.8-27b-run13-open-20261106-2dc461f52ea4.json embedding.with_panel["open-run13"].minus_without.blocks.hi; paper/data/run13-readers-open-model/simulated-readers-open-model-qwen3.8-27b-run13-open-20261106-2dc461f52ea4.json embedding.with_panel["open-run13"].minus_without.quarters.estimate; paper/data/run13-readers-open-model/simulated-readers-open-model-qwen3.8-27b-run13-open-20261106-2dc461f52ea4.json embedding.with_panel["open-run13"].minus_without.quarters.lo; paper/data/run13-readers-open-model/simulated-readers-open-model-qwen3.8-27b-run13-open-20261106-2dc461f52ea4.json embedding.with_panel["open-run13"].minus_without.quarters.hi
\quad with the open model's panel, minus without & -0.0023 (-0.0073, 0.0030) & -0.0009 (-0.0041, 0.0022) \\
\midrule
\multicolumn{3}{@{}l}{\emph{Calls of the open model}} \\
% src: paper/data/run13-readers-open-model/simulated-readers-qwen3.8-27b-run13-open-20261106-2dc461f52ea4.json counts.calls_planned; paper/data/run13-readers-open-model/simulated-readers-qwen3.8-27b-run13-open-20261106-2dc461f52ea4.json counts.calls_ok
\quad planned calls of the panel and probes, and usable replies & 9,007 & 9,007 \\
% src: paper/data/run13-readers-open-model/simulated-readers-open-model-qwen3.8-27b-run13-open-20261106-2dc461f52ea4.json direct.counts.calls_planned; paper/data/run13-readers-open-model/simulated-readers-open-model-qwen3.8-27b-run13-open-20261106-2dc461f52ea4.json direct.counts.calls_ok
\quad calls of the direct judgment, and usable replies & 616 & 616 \\
\bottomrule
\end{tabular}

\end{table*}

\begin{table*}[t]\centering\footnotesize
% src: tables/fmt-readers-probes.tex; src/mootcourt/simulators/readers/protocol.py N_PROBE_UNITS 300, STEM_WORDS 4, MIN_COMPLETION_WORDS 8
\caption{Three probes for earlier exposure to the archive, on 300 headline-only tests drawn with a fixed seed before any call. Completion: the model is given the first four words of a headline and asked for the rest; the measure is the token F1 of its completion, for the headline with the highest click rate of a test and for one other headline of that test, and the mean paired difference. A headline of fewer than eight words is skipped. Familiarity: the model rates how sure it is that it has seen each headline, and the rule that picks the higher rating is scored for direction. Named publisher: the model is told the publisher and the period, shown the two headlines of a contrast in random order, and asked which had the higher click rate. This third probe is a direct question and not a reader. Contrasts with equal ratings stay in the familiarity share of headlines that won, where each counts as half right, and are left out of the share that compares the ratings with the plain panel. The reports give their count.}
\label{tab:fmt-readers-probes}
% generated by paper/build_tables.py from the report keys named above each row (readers); it reads paper/data/run11-simulated-readers and paper/data/run12-simulated-readers-api
\begin{tabular}{p{7.6cm}rrr}
\toprule
probe and measure & Opus & Sonnet & Haiku \\
\midrule
\multicolumn{4}{@{}l}{\emph{Run 11, through the command line tool}} \\
% src: paper/data/run11-simulated-readers/simulated-readers-claude-opus-5-5-20261106-2dc461f52ea4.json probes.units; paper/data/run11-simulated-readers/simulated-readers-claude-sonnet-5-5-20261106-2dc461f52ea4.json probes.units; paper/data/run11-simulated-readers/simulated-readers-claude-haiku-5-5-20261106-2dc461f52ea4.json probes.units
tests probed & 300 & 300 & 300 \\
% src: paper/data/run11-simulated-readers/simulated-readers-claude-opus-5-5-20261106-2dc461f52ea4.json probes.completion.n_pairs; paper/data/run11-simulated-readers/simulated-readers-claude-sonnet-5-5-20261106-2dc461f52ea4.json probes.completion.n_pairs; paper/data/run11-simulated-readers/simulated-readers-claude-haiku-5-5-20261106-2dc461f52ea4.json probes.completion.n_pairs
completion: pairs of headlines scored & 295 & 295 & 295 \\
% src: paper/data/run11-simulated-readers/simulated-readers-claude-opus-5-5-20261106-2dc461f52ea4.json probes.completion.f1_winner; paper/data/run11-simulated-readers/simulated-readers-claude-sonnet-5-5-20261106-2dc461f52ea4.json probes.completion.f1_winner; paper/data/run11-simulated-readers/simulated-readers-claude-haiku-5-5-20261106-2dc461f52ea4.json probes.completion.f1_winner
completion: token F1, the winning headline & 0.155 & 0.136 & 0.102 \\
% src: paper/data/run11-simulated-readers/simulated-readers-claude-opus-5-5-20261106-2dc461f52ea4.json probes.completion.f1_other; paper/data/run11-simulated-readers/simulated-readers-claude-sonnet-5-5-20261106-2dc461f52ea4.json probes.completion.f1_other; paper/data/run11-simulated-readers/simulated-readers-claude-haiku-5-5-20261106-2dc461f52ea4.json probes.completion.f1_other
completion: token F1, another headline of the test & 0.145 & 0.145 & 0.106 \\
% src: paper/data/run11-simulated-readers/simulated-readers-claude-opus-5-5-20261106-2dc461f52ea4.json probes.completion.f1_difference; paper/data/run11-simulated-readers/simulated-readers-claude-opus-5-5-20261106-2dc461f52ea4.json probes.completion.f1_difference_se; paper/data/run11-simulated-readers/simulated-readers-claude-sonnet-5-5-20261106-2dc461f52ea4.json probes.completion.f1_difference; paper/data/run11-simulated-readers/simulated-readers-claude-sonnet-5-5-20261106-2dc461f52ea4.json probes.completion.f1_difference_se; paper/data/run11-simulated-readers/simulated-readers-claude-haiku-5-5-20261106-2dc461f52ea4.json probes.completion.f1_difference; paper/data/run11-simulated-readers/simulated-readers-claude-haiku-5-5-20261106-2dc461f52ea4.json probes.completion.f1_difference_se
completion: winner minus other (standard error) & 0.010 (0.008) & -0.009 (0.008) & -0.003 (0.008) \\
% src: paper/data/run11-simulated-readers/simulated-readers-claude-opus-5-5-20261106-2dc461f52ea4.json probes.familiarity.n; paper/data/run11-simulated-readers/simulated-readers-claude-sonnet-5-5-20261106-2dc461f52ea4.json probes.familiarity.n; paper/data/run11-simulated-readers/simulated-readers-claude-haiku-5-5-20261106-2dc461f52ea4.json probes.familiarity.n
familiarity: contrasts scored & 985 & 985 & 985 \\
% src: paper/data/run11-simulated-readers/simulated-readers-claude-opus-5-5-20261106-2dc461f52ea4.json probes.familiarity.direction_accuracy; paper/data/run11-simulated-readers/simulated-readers-claude-sonnet-5-5-20261106-2dc461f52ea4.json probes.familiarity.direction_accuracy; paper/data/run11-simulated-readers/simulated-readers-claude-haiku-5-5-20261106-2dc461f52ea4.json probes.familiarity.direction_accuracy
familiarity: the headline rated more familiar won, share & 0.526 & 0.534 & 0.525 \\
% src: paper/data/run11-simulated-readers/simulated-readers-claude-opus-5-5-20261106-2dc461f52ea4.json probes.familiarity.agreement_with_panel; paper/data/run11-simulated-readers/simulated-readers-claude-sonnet-5-5-20261106-2dc461f52ea4.json probes.familiarity.agreement_with_panel; paper/data/run11-simulated-readers/simulated-readers-claude-haiku-5-5-20261106-2dc461f52ea4.json probes.familiarity.agreement_with_panel
familiarity: same direction as the plain panel, share & 0.556 & 0.524 & 0.497 \\
% src: paper/data/run11-simulated-readers/simulated-readers-claude-opus-5-5-20261106-2dc461f52ea4.json probes.pairs.direction_accuracy; paper/data/run11-simulated-readers/simulated-readers-claude-sonnet-5-5-20261106-2dc461f52ea4.json probes.pairs.direction_accuracy; paper/data/run11-simulated-readers/simulated-readers-claude-haiku-5-5-20261106-2dc461f52ea4.json probes.pairs.direction_accuracy
named publisher: which headline of a pair won, share right & 0.595 & 0.557 & 0.557 \\
\midrule
\multicolumn{4}{@{}l}{\emph{Run 12, through the vendor's API}} \\
% src: paper/data/run12-simulated-readers-api/simulated-readers-claude-opus-5-5-run12-api-20261106-2dc461f52ea4.json probes.units; paper/data/run12-simulated-readers-api/simulated-readers-claude-sonnet-5-5-run12-api-20261106-2dc461f52ea4.json probes.units; paper/data/run12-simulated-readers-api/simulated-readers-claude-haiku-5-5-run12-api-20261106-2dc461f52ea4.json probes.units
tests probed & 300 & 300 & 300 \\
% src: paper/data/run12-simulated-readers-api/simulated-readers-claude-opus-5-5-run12-api-20261106-2dc461f52ea4.json probes.completion.n_pairs; paper/data/run12-simulated-readers-api/simulated-readers-claude-sonnet-5-5-run12-api-20261106-2dc461f52ea4.json probes.completion.n_pairs; paper/data/run12-simulated-readers-api/simulated-readers-claude-haiku-5-5-run12-api-20261106-2dc461f52ea4.json probes.completion.n_pairs
completion: pairs of headlines scored & 295 & 295 & 295 \\
% src: paper/data/run12-simulated-readers-api/simulated-readers-claude-opus-5-5-run12-api-20261106-2dc461f52ea4.json probes.completion.f1_winner; paper/data/run12-simulated-readers-api/simulated-readers-claude-sonnet-5-5-run12-api-20261106-2dc461f52ea4.json probes.completion.f1_winner; paper/data/run12-simulated-readers-api/simulated-readers-claude-haiku-5-5-run12-api-20261106-2dc461f52ea4.json probes.completion.f1_winner
completion: token F1, the winning headline & 0.149 & 0.135 & 0.109 \\
% src: paper/data/run12-simulated-readers-api/simulated-readers-claude-opus-5-5-run12-api-20261106-2dc461f52ea4.json probes.completion.f1_other; paper/data/run12-simulated-readers-api/simulated-readers-claude-sonnet-5-5-run12-api-20261106-2dc461f52ea4.json probes.completion.f1_other; paper/data/run12-simulated-readers-api/simulated-readers-claude-haiku-5-5-run12-api-20261106-2dc461f52ea4.json probes.completion.f1_other
completion: token F1, another headline of the test & 0.151 & 0.141 & 0.117 \\
% src: paper/data/run12-simulated-readers-api/simulated-readers-claude-opus-5-5-run12-api-20261106-2dc461f52ea4.json probes.completion.f1_difference; paper/data/run12-simulated-readers-api/simulated-readers-claude-opus-5-5-run12-api-20261106-2dc461f52ea4.json probes.completion.f1_difference_se; paper/data/run12-simulated-readers-api/simulated-readers-claude-sonnet-5-5-run12-api-20261106-2dc461f52ea4.json probes.completion.f1_difference; paper/data/run12-simulated-readers-api/simulated-readers-claude-sonnet-5-5-run12-api-20261106-2dc461f52ea4.json probes.completion.f1_difference_se; paper/data/run12-simulated-readers-api/simulated-readers-claude-haiku-5-5-run12-api-20261106-2dc461f52ea4.json probes.completion.f1_difference; paper/data/run12-simulated-readers-api/simulated-readers-claude-haiku-5-5-run12-api-20261106-2dc461f52ea4.json probes.completion.f1_difference_se
completion: winner minus other (standard error) & -0.001 (0.009) & -0.005 (0.008) & -0.007 (0.008) \\
% src: paper/data/run12-simulated-readers-api/simulated-readers-claude-opus-5-5-run12-api-20261106-2dc461f52ea4.json probes.familiarity.n; paper/data/run12-simulated-readers-api/simulated-readers-claude-sonnet-5-5-run12-api-20261106-2dc461f52ea4.json probes.familiarity.n; paper/data/run12-simulated-readers-api/simulated-readers-claude-haiku-5-5-run12-api-20261106-2dc461f52ea4.json probes.familiarity.n
familiarity: contrasts scored & 985 & 985 & 985 \\
% src: paper/data/run12-simulated-readers-api/simulated-readers-claude-opus-5-5-run12-api-20261106-2dc461f52ea4.json probes.familiarity.direction_accuracy; paper/data/run12-simulated-readers-api/simulated-readers-claude-sonnet-5-5-run12-api-20261106-2dc461f52ea4.json probes.familiarity.direction_accuracy; paper/data/run12-simulated-readers-api/simulated-readers-claude-haiku-5-5-run12-api-20261106-2dc461f52ea4.json probes.familiarity.direction_accuracy
familiarity: the headline rated more familiar won, share & 0.532 & 0.525 & 0.509 \\
% src: paper/data/run12-simulated-readers-api/simulated-readers-claude-opus-5-5-run12-api-20261106-2dc461f52ea4.json probes.familiarity.agreement_with_panel; paper/data/run12-simulated-readers-api/simulated-readers-claude-sonnet-5-5-run12-api-20261106-2dc461f52ea4.json probes.familiarity.agreement_with_panel; paper/data/run12-simulated-readers-api/simulated-readers-claude-haiku-5-5-run12-api-20261106-2dc461f52ea4.json probes.familiarity.agreement_with_panel
familiarity: same direction as the plain panel, share & 0.559 & 0.571 & 0.531 \\
% src: paper/data/run12-simulated-readers-api/simulated-readers-claude-opus-5-5-run12-api-20261106-2dc461f52ea4.json probes.pairs.direction_accuracy; paper/data/run12-simulated-readers-api/simulated-readers-claude-sonnet-5-5-run12-api-20261106-2dc461f52ea4.json probes.pairs.direction_accuracy; paper/data/run12-simulated-readers-api/simulated-readers-claude-haiku-5-5-run12-api-20261106-2dc461f52ea4.json probes.pairs.direction_accuracy
named publisher: which headline of a pair won, share right & 0.591 & 0.575 & 0.563 \\
\bottomrule
\end{tabular}

\end{table*}

\begin{table*}[t]\centering\footnotesize
% src: tables/fmt-readers-methods-api.tex; paper/data/run12-simulated-readers-api/simulated-readers-claude-opus-5-5-run12-api-20261106-2dc461f52ea4.json methods.audience_brief.minus_always_reference.resamples 500, counts.blocks 63; run 11 and run 12 are the names of the two run folders
\caption{The plain panel and the seven alignment methods in run 12, the same requests sent through the vendor's programming interface, by model. The first column is direction accuracy, and the other two are paired differences from the plain panel and from the earliest-arm rule, the rule that always picks the headline written first. Columns and intervals are those of Table~\ref{tab:fmt-readers-methods}, and no interval is adjusted for the number of comparisons.}
\label{tab:fmt-readers-methods-api}
% generated by paper/build_tables.py from the report keys named above each row (readers); it reads paper/data/run12-simulated-readers-api
\begin{tabular}{p{4.6cm}rrr}
\toprule
model and method & share matched & minus the plain panel & minus the earliest-arm rule \\
\midrule
\multicolumn{4}{@{}l}{\emph{Opus}} \\
% src: paper/data/run12-simulated-readers-api/simulated-readers-claude-opus-5-5-run12-api-20261106-2dc461f52ea4.json plain.accuracy; paper/data/run12-simulated-readers-api/simulated-readers-claude-opus-5-5-run12-api-20261106-2dc461f52ea4.json readings.plain_minus_always_reference.estimate; paper/data/run12-simulated-readers-api/simulated-readers-claude-opus-5-5-run12-api-20261106-2dc461f52ea4.json readings.plain_minus_always_reference.lo; paper/data/run12-simulated-readers-api/simulated-readers-claude-opus-5-5-run12-api-20261106-2dc461f52ea4.json readings.plain_minus_always_reference.hi
\quad plain panel & 0.5742 &  & 0.0041 (-0.0115, 0.0199) \\
% src: paper/data/run12-simulated-readers-api/simulated-readers-claude-opus-5-5-run12-api-20261106-2dc461f52ea4.json methods.direction_map.accuracy; paper/data/run12-simulated-readers-api/simulated-readers-claude-opus-5-5-run12-api-20261106-2dc461f52ea4.json methods.direction_map.minus_plain.estimate; paper/data/run12-simulated-readers-api/simulated-readers-claude-opus-5-5-run12-api-20261106-2dc461f52ea4.json methods.direction_map.minus_plain.lo; paper/data/run12-simulated-readers-api/simulated-readers-claude-opus-5-5-run12-api-20261106-2dc461f52ea4.json methods.direction_map.minus_plain.hi; paper/data/run12-simulated-readers-api/simulated-readers-claude-opus-5-5-run12-api-20261106-2dc461f52ea4.json methods.direction_map.minus_always_reference.estimate; paper/data/run12-simulated-readers-api/simulated-readers-claude-opus-5-5-run12-api-20261106-2dc461f52ea4.json methods.direction_map.minus_always_reference.lo; paper/data/run12-simulated-readers-api/simulated-readers-claude-opus-5-5-run12-api-20261106-2dc461f52ea4.json methods.direction_map.minus_always_reference.hi
\quad 1. direction map & 0.5855 & 0.0113 (-0.0016, 0.0238) & 0.0154 (0.0054, 0.0258) \\
% src: paper/data/run12-simulated-readers-api/simulated-readers-claude-opus-5-5-run12-api-20261106-2dc461f52ea4.json methods.reader_weights.accuracy; paper/data/run12-simulated-readers-api/simulated-readers-claude-opus-5-5-run12-api-20261106-2dc461f52ea4.json methods.reader_weights.minus_plain.estimate; paper/data/run12-simulated-readers-api/simulated-readers-claude-opus-5-5-run12-api-20261106-2dc461f52ea4.json methods.reader_weights.minus_plain.lo; paper/data/run12-simulated-readers-api/simulated-readers-claude-opus-5-5-run12-api-20261106-2dc461f52ea4.json methods.reader_weights.minus_plain.hi; paper/data/run12-simulated-readers-api/simulated-readers-claude-opus-5-5-run12-api-20261106-2dc461f52ea4.json methods.reader_weights.minus_always_reference.estimate; paper/data/run12-simulated-readers-api/simulated-readers-claude-opus-5-5-run12-api-20261106-2dc461f52ea4.json methods.reader_weights.minus_always_reference.lo; paper/data/run12-simulated-readers-api/simulated-readers-claude-opus-5-5-run12-api-20261106-2dc461f52ea4.json methods.reader_weights.minus_always_reference.hi
\quad 2. reader weights & 0.5851 & 0.0110 (-0.0003, 0.0218) & 0.0151 (-0.0036, 0.0326) \\
% src: paper/data/run12-simulated-readers-api/simulated-readers-claude-opus-5-5-run12-api-20261106-2dc461f52ea4.json methods.audience_matching.accuracy; paper/data/run12-simulated-readers-api/simulated-readers-claude-opus-5-5-run12-api-20261106-2dc461f52ea4.json methods.audience_matching.minus_plain.estimate; paper/data/run12-simulated-readers-api/simulated-readers-claude-opus-5-5-run12-api-20261106-2dc461f52ea4.json methods.audience_matching.minus_plain.lo; paper/data/run12-simulated-readers-api/simulated-readers-claude-opus-5-5-run12-api-20261106-2dc461f52ea4.json methods.audience_matching.minus_plain.hi; paper/data/run12-simulated-readers-api/simulated-readers-claude-opus-5-5-run12-api-20261106-2dc461f52ea4.json methods.audience_matching.minus_always_reference.estimate; paper/data/run12-simulated-readers-api/simulated-readers-claude-opus-5-5-run12-api-20261106-2dc461f52ea4.json methods.audience_matching.minus_always_reference.lo; paper/data/run12-simulated-readers-api/simulated-readers-claude-opus-5-5-run12-api-20261106-2dc461f52ea4.json methods.audience_matching.minus_always_reference.hi
\quad 3. audience matching & 0.5650 & -0.0091 (-0.0166, -0.0010) & -0.0050 (-0.0192, 0.0115) \\
% src: paper/data/run12-simulated-readers-api/simulated-readers-claude-opus-5-5-run12-api-20261106-2dc461f52ea4.json methods.error_correction.accuracy; paper/data/run12-simulated-readers-api/simulated-readers-claude-opus-5-5-run12-api-20261106-2dc461f52ea4.json methods.error_correction.minus_plain.estimate; paper/data/run12-simulated-readers-api/simulated-readers-claude-opus-5-5-run12-api-20261106-2dc461f52ea4.json methods.error_correction.minus_plain.lo; paper/data/run12-simulated-readers-api/simulated-readers-claude-opus-5-5-run12-api-20261106-2dc461f52ea4.json methods.error_correction.minus_plain.hi; paper/data/run12-simulated-readers-api/simulated-readers-claude-opus-5-5-run12-api-20261106-2dc461f52ea4.json methods.error_correction.minus_always_reference.estimate; paper/data/run12-simulated-readers-api/simulated-readers-claude-opus-5-5-run12-api-20261106-2dc461f52ea4.json methods.error_correction.minus_always_reference.lo; paper/data/run12-simulated-readers-api/simulated-readers-claude-opus-5-5-run12-api-20261106-2dc461f52ea4.json methods.error_correction.minus_always_reference.hi
\quad 4. error correction & 0.5805 & 0.0064 (-0.0020, 0.0145) & 0.0105 (-0.0056, 0.0266) \\
% src: paper/data/run12-simulated-readers-api/simulated-readers-claude-opus-5-5-run12-api-20261106-2dc461f52ea4.json methods.text_with_panel.accuracy; paper/data/run12-simulated-readers-api/simulated-readers-claude-opus-5-5-run12-api-20261106-2dc461f52ea4.json methods.text_with_panel.minus_plain.estimate; paper/data/run12-simulated-readers-api/simulated-readers-claude-opus-5-5-run12-api-20261106-2dc461f52ea4.json methods.text_with_panel.minus_plain.lo; paper/data/run12-simulated-readers-api/simulated-readers-claude-opus-5-5-run12-api-20261106-2dc461f52ea4.json methods.text_with_panel.minus_plain.hi; paper/data/run12-simulated-readers-api/simulated-readers-claude-opus-5-5-run12-api-20261106-2dc461f52ea4.json methods.text_with_panel.minus_always_reference.estimate; paper/data/run12-simulated-readers-api/simulated-readers-claude-opus-5-5-run12-api-20261106-2dc461f52ea4.json methods.text_with_panel.minus_always_reference.lo; paper/data/run12-simulated-readers-api/simulated-readers-claude-opus-5-5-run12-api-20261106-2dc461f52ea4.json methods.text_with_panel.minus_always_reference.hi
\quad 5. text classifier with the panel & 0.5669 & -0.0073 (-0.0231, 0.0069) & -0.0032 (-0.0182, 0.0111) \\
% src: paper/data/run12-simulated-readers-api/simulated-readers-claude-opus-5-5-run12-api-20261106-2dc461f52ea4.json methods.examples_in_prompt.accuracy; paper/data/run12-simulated-readers-api/simulated-readers-claude-opus-5-5-run12-api-20261106-2dc461f52ea4.json methods.examples_in_prompt.minus_plain.estimate; paper/data/run12-simulated-readers-api/simulated-readers-claude-opus-5-5-run12-api-20261106-2dc461f52ea4.json methods.examples_in_prompt.minus_plain.lo; paper/data/run12-simulated-readers-api/simulated-readers-claude-opus-5-5-run12-api-20261106-2dc461f52ea4.json methods.examples_in_prompt.minus_plain.hi; paper/data/run12-simulated-readers-api/simulated-readers-claude-opus-5-5-run12-api-20261106-2dc461f52ea4.json methods.examples_in_prompt.minus_always_reference.estimate; paper/data/run12-simulated-readers-api/simulated-readers-claude-opus-5-5-run12-api-20261106-2dc461f52ea4.json methods.examples_in_prompt.minus_always_reference.lo; paper/data/run12-simulated-readers-api/simulated-readers-claude-opus-5-5-run12-api-20261106-2dc461f52ea4.json methods.examples_in_prompt.minus_always_reference.hi
\quad 6. earlier results in the prompt & 0.5880 & 0.0138 (0.0043, 0.0244) & 0.0179 (0.0026, 0.0328) \\
% src: paper/data/run12-simulated-readers-api/simulated-readers-claude-opus-5-5-run12-api-20261106-2dc461f52ea4.json methods.audience_brief.accuracy; paper/data/run12-simulated-readers-api/simulated-readers-claude-opus-5-5-run12-api-20261106-2dc461f52ea4.json methods.audience_brief.minus_plain.estimate; paper/data/run12-simulated-readers-api/simulated-readers-claude-opus-5-5-run12-api-20261106-2dc461f52ea4.json methods.audience_brief.minus_plain.lo; paper/data/run12-simulated-readers-api/simulated-readers-claude-opus-5-5-run12-api-20261106-2dc461f52ea4.json methods.audience_brief.minus_plain.hi; paper/data/run12-simulated-readers-api/simulated-readers-claude-opus-5-5-run12-api-20261106-2dc461f52ea4.json methods.audience_brief.minus_always_reference.estimate; paper/data/run12-simulated-readers-api/simulated-readers-claude-opus-5-5-run12-api-20261106-2dc461f52ea4.json methods.audience_brief.minus_always_reference.lo; paper/data/run12-simulated-readers-api/simulated-readers-claude-opus-5-5-run12-api-20261106-2dc461f52ea4.json methods.audience_brief.minus_always_reference.hi
\quad 7. audience brief & 0.6170 & 0.0428 (0.0292, 0.0542) & 0.0469 (0.0300, 0.0623) \\
\midrule
\multicolumn{4}{@{}l}{\emph{Sonnet}} \\
% src: paper/data/run12-simulated-readers-api/simulated-readers-claude-sonnet-5-5-run12-api-20261106-2dc461f52ea4.json plain.accuracy; paper/data/run12-simulated-readers-api/simulated-readers-claude-sonnet-5-5-run12-api-20261106-2dc461f52ea4.json readings.plain_minus_always_reference.estimate; paper/data/run12-simulated-readers-api/simulated-readers-claude-sonnet-5-5-run12-api-20261106-2dc461f52ea4.json readings.plain_minus_always_reference.lo; paper/data/run12-simulated-readers-api/simulated-readers-claude-sonnet-5-5-run12-api-20261106-2dc461f52ea4.json readings.plain_minus_always_reference.hi
\quad plain panel & 0.5279 &  & -0.0422 (-0.0582, -0.0252) \\
% src: paper/data/run12-simulated-readers-api/simulated-readers-claude-sonnet-5-5-run12-api-20261106-2dc461f52ea4.json methods.direction_map.accuracy; paper/data/run12-simulated-readers-api/simulated-readers-claude-sonnet-5-5-run12-api-20261106-2dc461f52ea4.json methods.direction_map.minus_plain.estimate; paper/data/run12-simulated-readers-api/simulated-readers-claude-sonnet-5-5-run12-api-20261106-2dc461f52ea4.json methods.direction_map.minus_plain.lo; paper/data/run12-simulated-readers-api/simulated-readers-claude-sonnet-5-5-run12-api-20261106-2dc461f52ea4.json methods.direction_map.minus_plain.hi; paper/data/run12-simulated-readers-api/simulated-readers-claude-sonnet-5-5-run12-api-20261106-2dc461f52ea4.json methods.direction_map.minus_always_reference.estimate; paper/data/run12-simulated-readers-api/simulated-readers-claude-sonnet-5-5-run12-api-20261106-2dc461f52ea4.json methods.direction_map.minus_always_reference.lo; paper/data/run12-simulated-readers-api/simulated-readers-claude-sonnet-5-5-run12-api-20261106-2dc461f52ea4.json methods.direction_map.minus_always_reference.hi
\quad 1. direction map & 0.5670 & 0.0391 (0.0233, 0.0531) & -0.0030 (-0.0094, 0.0024) \\
% src: paper/data/run12-simulated-readers-api/simulated-readers-claude-sonnet-5-5-run12-api-20261106-2dc461f52ea4.json methods.reader_weights.accuracy; paper/data/run12-simulated-readers-api/simulated-readers-claude-sonnet-5-5-run12-api-20261106-2dc461f52ea4.json methods.reader_weights.minus_plain.estimate; paper/data/run12-simulated-readers-api/simulated-readers-claude-sonnet-5-5-run12-api-20261106-2dc461f52ea4.json methods.reader_weights.minus_plain.lo; paper/data/run12-simulated-readers-api/simulated-readers-claude-sonnet-5-5-run12-api-20261106-2dc461f52ea4.json methods.reader_weights.minus_plain.hi; paper/data/run12-simulated-readers-api/simulated-readers-claude-sonnet-5-5-run12-api-20261106-2dc461f52ea4.json methods.reader_weights.minus_always_reference.estimate; paper/data/run12-simulated-readers-api/simulated-readers-claude-sonnet-5-5-run12-api-20261106-2dc461f52ea4.json methods.reader_weights.minus_always_reference.lo; paper/data/run12-simulated-readers-api/simulated-readers-claude-sonnet-5-5-run12-api-20261106-2dc461f52ea4.json methods.reader_weights.minus_always_reference.hi
\quad 2. reader weights & 0.5605 & 0.0326 (0.0196, 0.0448) & -0.0096 (-0.0283, 0.0091) \\
% src: paper/data/run12-simulated-readers-api/simulated-readers-claude-sonnet-5-5-run12-api-20261106-2dc461f52ea4.json methods.audience_matching.accuracy; paper/data/run12-simulated-readers-api/simulated-readers-claude-sonnet-5-5-run12-api-20261106-2dc461f52ea4.json methods.audience_matching.minus_plain.estimate; paper/data/run12-simulated-readers-api/simulated-readers-claude-sonnet-5-5-run12-api-20261106-2dc461f52ea4.json methods.audience_matching.minus_plain.lo; paper/data/run12-simulated-readers-api/simulated-readers-claude-sonnet-5-5-run12-api-20261106-2dc461f52ea4.json methods.audience_matching.minus_plain.hi; paper/data/run12-simulated-readers-api/simulated-readers-claude-sonnet-5-5-run12-api-20261106-2dc461f52ea4.json methods.audience_matching.minus_always_reference.estimate; paper/data/run12-simulated-readers-api/simulated-readers-claude-sonnet-5-5-run12-api-20261106-2dc461f52ea4.json methods.audience_matching.minus_always_reference.lo; paper/data/run12-simulated-readers-api/simulated-readers-claude-sonnet-5-5-run12-api-20261106-2dc461f52ea4.json methods.audience_matching.minus_always_reference.hi
\quad 3. audience matching & 0.5283 & 0.0004 (-0.0106, 0.0101) & -0.0417 (-0.0614, -0.0250) \\
% src: paper/data/run12-simulated-readers-api/simulated-readers-claude-sonnet-5-5-run12-api-20261106-2dc461f52ea4.json methods.error_correction.accuracy; paper/data/run12-simulated-readers-api/simulated-readers-claude-sonnet-5-5-run12-api-20261106-2dc461f52ea4.json methods.error_correction.minus_plain.estimate; paper/data/run12-simulated-readers-api/simulated-readers-claude-sonnet-5-5-run12-api-20261106-2dc461f52ea4.json methods.error_correction.minus_plain.lo; paper/data/run12-simulated-readers-api/simulated-readers-claude-sonnet-5-5-run12-api-20261106-2dc461f52ea4.json methods.error_correction.minus_plain.hi; paper/data/run12-simulated-readers-api/simulated-readers-claude-sonnet-5-5-run12-api-20261106-2dc461f52ea4.json methods.error_correction.minus_always_reference.estimate; paper/data/run12-simulated-readers-api/simulated-readers-claude-sonnet-5-5-run12-api-20261106-2dc461f52ea4.json methods.error_correction.minus_always_reference.lo; paper/data/run12-simulated-readers-api/simulated-readers-claude-sonnet-5-5-run12-api-20261106-2dc461f52ea4.json methods.error_correction.minus_always_reference.hi
\quad 4. error correction & 0.5458 & 0.0179 (0.0073, 0.0262) & -0.0243 (-0.0417, -0.0075) \\
% src: paper/data/run12-simulated-readers-api/simulated-readers-claude-sonnet-5-5-run12-api-20261106-2dc461f52ea4.json methods.text_with_panel.accuracy; paper/data/run12-simulated-readers-api/simulated-readers-claude-sonnet-5-5-run12-api-20261106-2dc461f52ea4.json methods.text_with_panel.minus_plain.estimate; paper/data/run12-simulated-readers-api/simulated-readers-claude-sonnet-5-5-run12-api-20261106-2dc461f52ea4.json methods.text_with_panel.minus_plain.lo; paper/data/run12-simulated-readers-api/simulated-readers-claude-sonnet-5-5-run12-api-20261106-2dc461f52ea4.json methods.text_with_panel.minus_plain.hi; paper/data/run12-simulated-readers-api/simulated-readers-claude-sonnet-5-5-run12-api-20261106-2dc461f52ea4.json methods.text_with_panel.minus_always_reference.estimate; paper/data/run12-simulated-readers-api/simulated-readers-claude-sonnet-5-5-run12-api-20261106-2dc461f52ea4.json methods.text_with_panel.minus_always_reference.lo; paper/data/run12-simulated-readers-api/simulated-readers-claude-sonnet-5-5-run12-api-20261106-2dc461f52ea4.json methods.text_with_panel.minus_always_reference.hi
\quad 5. text classifier with the panel & 0.5478 & 0.0199 (0.0033, 0.0346) & -0.0223 (-0.0386, -0.0057) \\
% src: paper/data/run12-simulated-readers-api/simulated-readers-claude-sonnet-5-5-run12-api-20261106-2dc461f52ea4.json methods.examples_in_prompt.accuracy; paper/data/run12-simulated-readers-api/simulated-readers-claude-sonnet-5-5-run12-api-20261106-2dc461f52ea4.json methods.examples_in_prompt.minus_plain.estimate; paper/data/run12-simulated-readers-api/simulated-readers-claude-sonnet-5-5-run12-api-20261106-2dc461f52ea4.json methods.examples_in_prompt.minus_plain.lo; paper/data/run12-simulated-readers-api/simulated-readers-claude-sonnet-5-5-run12-api-20261106-2dc461f52ea4.json methods.examples_in_prompt.minus_plain.hi; paper/data/run12-simulated-readers-api/simulated-readers-claude-sonnet-5-5-run12-api-20261106-2dc461f52ea4.json methods.examples_in_prompt.minus_always_reference.estimate; paper/data/run12-simulated-readers-api/simulated-readers-claude-sonnet-5-5-run12-api-20261106-2dc461f52ea4.json methods.examples_in_prompt.minus_always_reference.lo; paper/data/run12-simulated-readers-api/simulated-readers-claude-sonnet-5-5-run12-api-20261106-2dc461f52ea4.json methods.examples_in_prompt.minus_always_reference.hi
\quad 6. earlier results in the prompt & 0.5462 & 0.0183 (0.0077, 0.0273) & -0.0239 (-0.0418, -0.0063) \\
% src: paper/data/run12-simulated-readers-api/simulated-readers-claude-sonnet-5-5-run12-api-20261106-2dc461f52ea4.json methods.audience_brief.accuracy; paper/data/run12-simulated-readers-api/simulated-readers-claude-sonnet-5-5-run12-api-20261106-2dc461f52ea4.json methods.audience_brief.minus_plain.estimate; paper/data/run12-simulated-readers-api/simulated-readers-claude-sonnet-5-5-run12-api-20261106-2dc461f52ea4.json methods.audience_brief.minus_plain.lo; paper/data/run12-simulated-readers-api/simulated-readers-claude-sonnet-5-5-run12-api-20261106-2dc461f52ea4.json methods.audience_brief.minus_plain.hi; paper/data/run12-simulated-readers-api/simulated-readers-claude-sonnet-5-5-run12-api-20261106-2dc461f52ea4.json methods.audience_brief.minus_always_reference.estimate; paper/data/run12-simulated-readers-api/simulated-readers-claude-sonnet-5-5-run12-api-20261106-2dc461f52ea4.json methods.audience_brief.minus_always_reference.lo; paper/data/run12-simulated-readers-api/simulated-readers-claude-sonnet-5-5-run12-api-20261106-2dc461f52ea4.json methods.audience_brief.minus_always_reference.hi
\quad 7. audience brief & 0.5878 & 0.0599 (0.0481, 0.0717) & 0.0178 (0.0010, 0.0352) \\
\midrule
\multicolumn{4}{@{}l}{\emph{Haiku}} \\
% src: paper/data/run12-simulated-readers-api/simulated-readers-claude-haiku-5-5-run12-api-20261106-2dc461f52ea4.json plain.accuracy; paper/data/run12-simulated-readers-api/simulated-readers-claude-haiku-5-5-run12-api-20261106-2dc461f52ea4.json readings.plain_minus_always_reference.estimate; paper/data/run12-simulated-readers-api/simulated-readers-claude-haiku-5-5-run12-api-20261106-2dc461f52ea4.json readings.plain_minus_always_reference.lo; paper/data/run12-simulated-readers-api/simulated-readers-claude-haiku-5-5-run12-api-20261106-2dc461f52ea4.json readings.plain_minus_always_reference.hi
\quad plain panel & 0.5223 &  & -0.0478 (-0.0666, -0.0308) \\
% src: paper/data/run12-simulated-readers-api/simulated-readers-claude-haiku-5-5-run12-api-20261106-2dc461f52ea4.json methods.direction_map.accuracy; paper/data/run12-simulated-readers-api/simulated-readers-claude-haiku-5-5-run12-api-20261106-2dc461f52ea4.json methods.direction_map.minus_plain.estimate; paper/data/run12-simulated-readers-api/simulated-readers-claude-haiku-5-5-run12-api-20261106-2dc461f52ea4.json methods.direction_map.minus_plain.lo; paper/data/run12-simulated-readers-api/simulated-readers-claude-haiku-5-5-run12-api-20261106-2dc461f52ea4.json methods.direction_map.minus_plain.hi; paper/data/run12-simulated-readers-api/simulated-readers-claude-haiku-5-5-run12-api-20261106-2dc461f52ea4.json methods.direction_map.minus_always_reference.estimate; paper/data/run12-simulated-readers-api/simulated-readers-claude-haiku-5-5-run12-api-20261106-2dc461f52ea4.json methods.direction_map.minus_always_reference.lo; paper/data/run12-simulated-readers-api/simulated-readers-claude-haiku-5-5-run12-api-20261106-2dc461f52ea4.json methods.direction_map.minus_always_reference.hi
\quad 1. direction map & 0.5608 & 0.0386 (0.0227, 0.0551) & -0.0092 (-0.0163, -0.0024) \\
% src: paper/data/run12-simulated-readers-api/simulated-readers-claude-haiku-5-5-run12-api-20261106-2dc461f52ea4.json methods.reader_weights.accuracy; paper/data/run12-simulated-readers-api/simulated-readers-claude-haiku-5-5-run12-api-20261106-2dc461f52ea4.json methods.reader_weights.minus_plain.estimate; paper/data/run12-simulated-readers-api/simulated-readers-claude-haiku-5-5-run12-api-20261106-2dc461f52ea4.json methods.reader_weights.minus_plain.lo; paper/data/run12-simulated-readers-api/simulated-readers-claude-haiku-5-5-run12-api-20261106-2dc461f52ea4.json methods.reader_weights.minus_plain.hi; paper/data/run12-simulated-readers-api/simulated-readers-claude-haiku-5-5-run12-api-20261106-2dc461f52ea4.json methods.reader_weights.minus_always_reference.estimate; paper/data/run12-simulated-readers-api/simulated-readers-claude-haiku-5-5-run12-api-20261106-2dc461f52ea4.json methods.reader_weights.minus_always_reference.lo; paper/data/run12-simulated-readers-api/simulated-readers-claude-haiku-5-5-run12-api-20261106-2dc461f52ea4.json methods.reader_weights.minus_always_reference.hi
\quad 2. reader weights & 0.5211 & -0.0012 (-0.0133, 0.0122) & -0.0489 (-0.0693, -0.0284) \\
% src: paper/data/run12-simulated-readers-api/simulated-readers-claude-haiku-5-5-run12-api-20261106-2dc461f52ea4.json methods.audience_matching.accuracy; paper/data/run12-simulated-readers-api/simulated-readers-claude-haiku-5-5-run12-api-20261106-2dc461f52ea4.json methods.audience_matching.minus_plain.estimate; paper/data/run12-simulated-readers-api/simulated-readers-claude-haiku-5-5-run12-api-20261106-2dc461f52ea4.json methods.audience_matching.minus_plain.lo; paper/data/run12-simulated-readers-api/simulated-readers-claude-haiku-5-5-run12-api-20261106-2dc461f52ea4.json methods.audience_matching.minus_plain.hi; paper/data/run12-simulated-readers-api/simulated-readers-claude-haiku-5-5-run12-api-20261106-2dc461f52ea4.json methods.audience_matching.minus_always_reference.estimate; paper/data/run12-simulated-readers-api/simulated-readers-claude-haiku-5-5-run12-api-20261106-2dc461f52ea4.json methods.audience_matching.minus_always_reference.lo; paper/data/run12-simulated-readers-api/simulated-readers-claude-haiku-5-5-run12-api-20261106-2dc461f52ea4.json methods.audience_matching.minus_always_reference.hi
\quad 3. audience matching & 0.5082 & -0.0141 (-0.0254, -0.0045) & -0.0619 (-0.0798, -0.0440) \\
% src: paper/data/run12-simulated-readers-api/simulated-readers-claude-haiku-5-5-run12-api-20261106-2dc461f52ea4.json methods.error_correction.accuracy; paper/data/run12-simulated-readers-api/simulated-readers-claude-haiku-5-5-run12-api-20261106-2dc461f52ea4.json methods.error_correction.minus_plain.estimate; paper/data/run12-simulated-readers-api/simulated-readers-claude-haiku-5-5-run12-api-20261106-2dc461f52ea4.json methods.error_correction.minus_plain.lo; paper/data/run12-simulated-readers-api/simulated-readers-claude-haiku-5-5-run12-api-20261106-2dc461f52ea4.json methods.error_correction.minus_plain.hi; paper/data/run12-simulated-readers-api/simulated-readers-claude-haiku-5-5-run12-api-20261106-2dc461f52ea4.json methods.error_correction.minus_always_reference.estimate; paper/data/run12-simulated-readers-api/simulated-readers-claude-haiku-5-5-run12-api-20261106-2dc461f52ea4.json methods.error_correction.minus_always_reference.lo; paper/data/run12-simulated-readers-api/simulated-readers-claude-haiku-5-5-run12-api-20261106-2dc461f52ea4.json methods.error_correction.minus_always_reference.hi
\quad 4. error correction & 0.5406 & 0.0183 (0.0066, 0.0290) & -0.0295 (-0.0484, -0.0124) \\
% src: paper/data/run12-simulated-readers-api/simulated-readers-claude-haiku-5-5-run12-api-20261106-2dc461f52ea4.json methods.text_with_panel.accuracy; paper/data/run12-simulated-readers-api/simulated-readers-claude-haiku-5-5-run12-api-20261106-2dc461f52ea4.json methods.text_with_panel.minus_plain.estimate; paper/data/run12-simulated-readers-api/simulated-readers-claude-haiku-5-5-run12-api-20261106-2dc461f52ea4.json methods.text_with_panel.minus_plain.lo; paper/data/run12-simulated-readers-api/simulated-readers-claude-haiku-5-5-run12-api-20261106-2dc461f52ea4.json methods.text_with_panel.minus_plain.hi; paper/data/run12-simulated-readers-api/simulated-readers-claude-haiku-5-5-run12-api-20261106-2dc461f52ea4.json methods.text_with_panel.minus_always_reference.estimate; paper/data/run12-simulated-readers-api/simulated-readers-claude-haiku-5-5-run12-api-20261106-2dc461f52ea4.json methods.text_with_panel.minus_always_reference.lo; paper/data/run12-simulated-readers-api/simulated-readers-claude-haiku-5-5-run12-api-20261106-2dc461f52ea4.json methods.text_with_panel.minus_always_reference.hi
\quad 5. text classifier with the panel & 0.5536 & 0.0313 (0.0138, 0.0497) & -0.0164 (-0.0318, 0.0012) \\
% src: paper/data/run12-simulated-readers-api/simulated-readers-claude-haiku-5-5-run12-api-20261106-2dc461f52ea4.json methods.examples_in_prompt.accuracy; paper/data/run12-simulated-readers-api/simulated-readers-claude-haiku-5-5-run12-api-20261106-2dc461f52ea4.json methods.examples_in_prompt.minus_plain.estimate; paper/data/run12-simulated-readers-api/simulated-readers-claude-haiku-5-5-run12-api-20261106-2dc461f52ea4.json methods.examples_in_prompt.minus_plain.lo; paper/data/run12-simulated-readers-api/simulated-readers-claude-haiku-5-5-run12-api-20261106-2dc461f52ea4.json methods.examples_in_prompt.minus_plain.hi; paper/data/run12-simulated-readers-api/simulated-readers-claude-haiku-5-5-run12-api-20261106-2dc461f52ea4.json methods.examples_in_prompt.minus_always_reference.estimate; paper/data/run12-simulated-readers-api/simulated-readers-claude-haiku-5-5-run12-api-20261106-2dc461f52ea4.json methods.examples_in_prompt.minus_always_reference.lo; paper/data/run12-simulated-readers-api/simulated-readers-claude-haiku-5-5-run12-api-20261106-2dc461f52ea4.json methods.examples_in_prompt.minus_always_reference.hi
\quad 6. earlier results in the prompt & 0.5227 & 0.0004 (-0.0119, 0.0123) & -0.0474 (-0.0673, -0.0284) \\
% src: paper/data/run12-simulated-readers-api/simulated-readers-claude-haiku-5-5-run12-api-20261106-2dc461f52ea4.json methods.audience_brief.accuracy; paper/data/run12-simulated-readers-api/simulated-readers-claude-haiku-5-5-run12-api-20261106-2dc461f52ea4.json methods.audience_brief.minus_plain.estimate; paper/data/run12-simulated-readers-api/simulated-readers-claude-haiku-5-5-run12-api-20261106-2dc461f52ea4.json methods.audience_brief.minus_plain.lo; paper/data/run12-simulated-readers-api/simulated-readers-claude-haiku-5-5-run12-api-20261106-2dc461f52ea4.json methods.audience_brief.minus_plain.hi; paper/data/run12-simulated-readers-api/simulated-readers-claude-haiku-5-5-run12-api-20261106-2dc461f52ea4.json methods.audience_brief.minus_always_reference.estimate; paper/data/run12-simulated-readers-api/simulated-readers-claude-haiku-5-5-run12-api-20261106-2dc461f52ea4.json methods.audience_brief.minus_always_reference.lo; paper/data/run12-simulated-readers-api/simulated-readers-claude-haiku-5-5-run12-api-20261106-2dc461f52ea4.json methods.audience_brief.minus_always_reference.hi
\quad 7. audience brief & 0.5401 & 0.0178 (0.0046, 0.0297) & -0.0300 (-0.0472, -0.0131) \\
\bottomrule
\end{tabular}

\end{table*}

\begin{table*}[t]\centering\footnotesize
% src: tables/fmt-readers-more.tex; paper/data/run11-simulated-readers/simulated-readers-claude-opus-5-5-20261106-2dc461f52ea4.json ranking.frac 0.25, ranking.resamples 500, plain.decisive.n 1634; paper/data/run9-visible-exposure/backtest-upworthy-20261106-2dc461f52ea4.json tables.aggregate[cell=headline_only,rung=always_reference].direction_accuracy_decisive (the row of the earliest-arm rule on decisive contrasts, 0.6481 against 0.5701 on all contrasts), tables.aggregate[cell=headline_only,rung=always_reference].n_decisive 1634; run 11 and run 12 are the names of the two run folders
\caption{Measures of run 11 that carry no verdict. Decisive contrasts are those whose observed effect exceeds 1.96 standard errors in size. On the decisive contrasts the earliest-arm rule is itself right more often than on all contrasts, so the plain panel's row there is read against that rule and not against one half. The rule's row comes from the main run's report, on the same contrasts. The two halves of the panel are readers 1 to 5 and readers 6 to 10. The Brier skill of the direction map is against the share of earlier contrasts won by the variant. This skill is for the direction of a contrast and not for the significant-positive event, so it is not comparable with the Brier skills of Section~\ref{sec:ladder-results}. A unit is a test at one origin, and one test is scored at two origins. The arm-count rule and the random order matched on arm count are those of Section~\ref{sec:protocol}. For ranking, a unit's score is the largest panel difference among its contrasts, and captured lift is the share of the total positive lift held by the one fourth of the units ranked first within each origin. Intervals are at 90 percent from 500 resamples of blocks and are not adjusted.}
\label{tab:fmt-readers-more}
% generated by paper/build_tables.py from the report keys named above each row (readers); it reads paper/data/run11-simulated-readers and, for one row, paper/data/run9-visible-exposure
\begin{tabular}{p{6.6cm}rrr}
\toprule
measure & Opus & Sonnet & Haiku \\
\midrule
% src: paper/data/run11-simulated-readers/simulated-readers-claude-opus-5-5-20261106-2dc461f52ea4.json plain.accuracy_minus_half.estimate; paper/data/run11-simulated-readers/simulated-readers-claude-opus-5-5-20261106-2dc461f52ea4.json plain.accuracy_minus_half.lo; paper/data/run11-simulated-readers/simulated-readers-claude-opus-5-5-20261106-2dc461f52ea4.json plain.accuracy_minus_half.hi; paper/data/run11-simulated-readers/simulated-readers-claude-sonnet-5-5-20261106-2dc461f52ea4.json plain.accuracy_minus_half.estimate; paper/data/run11-simulated-readers/simulated-readers-claude-sonnet-5-5-20261106-2dc461f52ea4.json plain.accuracy_minus_half.lo; paper/data/run11-simulated-readers/simulated-readers-claude-sonnet-5-5-20261106-2dc461f52ea4.json plain.accuracy_minus_half.hi; paper/data/run11-simulated-readers/simulated-readers-claude-haiku-5-5-20261106-2dc461f52ea4.json plain.accuracy_minus_half.estimate; paper/data/run11-simulated-readers/simulated-readers-claude-haiku-5-5-20261106-2dc461f52ea4.json plain.accuracy_minus_half.lo; paper/data/run11-simulated-readers/simulated-readers-claude-haiku-5-5-20261106-2dc461f52ea4.json plain.accuracy_minus_half.hi
plain panel minus one half & 0.0742 (0.0616, 0.0878) & 0.0234 (0.0118, 0.0362) & 0.0144 (0.0034, 0.0256) \\
% src: paper/data/run11-simulated-readers/simulated-readers-claude-opus-5-5-20261106-2dc461f52ea4.json plain.decisive.n; paper/data/run11-simulated-readers/simulated-readers-claude-sonnet-5-5-20261106-2dc461f52ea4.json plain.decisive.n; paper/data/run11-simulated-readers/simulated-readers-claude-haiku-5-5-20261106-2dc461f52ea4.json plain.decisive.n
decisive contrasts & 1,634 & 1,634 & 1,634 \\
% src: paper/data/run11-simulated-readers/simulated-readers-claude-opus-5-5-20261106-2dc461f52ea4.json plain.decisive.accuracy; paper/data/run11-simulated-readers/simulated-readers-claude-sonnet-5-5-20261106-2dc461f52ea4.json plain.decisive.accuracy; paper/data/run11-simulated-readers/simulated-readers-claude-haiku-5-5-20261106-2dc461f52ea4.json plain.decisive.accuracy
plain panel, share matched on decisive contrasts & 0.6392 & 0.5692 & 0.5459 \\
% src: paper/data/run9-visible-exposure/backtest-upworthy-20261106-2dc461f52ea4.json tables.aggregate[cell=headline_only,rung=always_reference].direction_accuracy_decisive; paper/data/run9-visible-exposure/backtest-upworthy-20261106-2dc461f52ea4.json tables.aggregate[cell=headline_only,rung=always_reference].direction_accuracy_decisive; paper/data/run9-visible-exposure/backtest-upworthy-20261106-2dc461f52ea4.json tables.aggregate[cell=headline_only,rung=always_reference].direction_accuracy_decisive
the earliest-arm rule, share matched on decisive contrasts & 0.6481 & 0.6481 & 0.6481 \\
% src: paper/data/run11-simulated-readers/simulated-readers-claude-opus-5-5-20261106-2dc461f52ea4.json plain.split_half.agree_share; paper/data/run11-simulated-readers/simulated-readers-claude-sonnet-5-5-20261106-2dc461f52ea4.json plain.split_half.agree_share; paper/data/run11-simulated-readers/simulated-readers-claude-haiku-5-5-20261106-2dc461f52ea4.json plain.split_half.agree_share
the two halves of the panel give the same direction, share & 0.750 & 0.746 & 0.660 \\
% src: paper/data/run11-simulated-readers/simulated-readers-claude-opus-5-5-20261106-2dc461f52ea4.json plain.ties; paper/data/run11-simulated-readers/simulated-readers-claude-sonnet-5-5-20261106-2dc461f52ea4.json plain.ties; paper/data/run11-simulated-readers/simulated-readers-claude-haiku-5-5-20261106-2dc461f52ea4.json plain.ties
contrasts on which the plain panel is tied & 3 & 1 & 0 \\
% src: paper/data/run11-simulated-readers/simulated-readers-claude-opus-5-5-20261106-2dc461f52ea4.json plain.no_reader; paper/data/run11-simulated-readers/simulated-readers-claude-sonnet-5-5-20261106-2dc461f52ea4.json plain.no_reader; paper/data/run11-simulated-readers/simulated-readers-claude-haiku-5-5-20261106-2dc461f52ea4.json plain.no_reader
contrasts left with no reader & 0 & 0 & 0 \\
% src: paper/data/run11-simulated-readers/simulated-readers-claude-opus-5-5-20261106-2dc461f52ea4.json methods.direction_map.brier_skill; paper/data/run11-simulated-readers/simulated-readers-claude-sonnet-5-5-20261106-2dc461f52ea4.json methods.direction_map.brier_skill; paper/data/run11-simulated-readers/simulated-readers-claude-haiku-5-5-20261106-2dc461f52ea4.json methods.direction_map.brier_skill
direction map, Brier skill & 0.0218 & 0.0029 & 0.0015 \\
% src: paper/data/run11-simulated-readers/simulated-readers-claude-opus-5-5-20261106-2dc461f52ea4.json baselines.text_lr_minus_always_reference.estimate; paper/data/run11-simulated-readers/simulated-readers-claude-opus-5-5-20261106-2dc461f52ea4.json baselines.text_lr_minus_always_reference.lo; paper/data/run11-simulated-readers/simulated-readers-claude-opus-5-5-20261106-2dc461f52ea4.json baselines.text_lr_minus_always_reference.hi; paper/data/run11-simulated-readers/simulated-readers-claude-sonnet-5-5-20261106-2dc461f52ea4.json baselines.text_lr_minus_always_reference.estimate; paper/data/run11-simulated-readers/simulated-readers-claude-sonnet-5-5-20261106-2dc461f52ea4.json baselines.text_lr_minus_always_reference.lo; paper/data/run11-simulated-readers/simulated-readers-claude-sonnet-5-5-20261106-2dc461f52ea4.json baselines.text_lr_minus_always_reference.hi; paper/data/run11-simulated-readers/simulated-readers-claude-haiku-5-5-20261106-2dc461f52ea4.json baselines.text_lr_minus_always_reference.estimate; paper/data/run11-simulated-readers/simulated-readers-claude-haiku-5-5-20261106-2dc461f52ea4.json baselines.text_lr_minus_always_reference.lo; paper/data/run11-simulated-readers/simulated-readers-claude-haiku-5-5-20261106-2dc461f52ea4.json baselines.text_lr_minus_always_reference.hi
text classifier minus the earliest-arm rule & -0.0216 (-0.0365, -0.0070) & -0.0216 (-0.0365, -0.0070) & -0.0216 (-0.0365, -0.0070) \\
\midrule
\multicolumn{4}{@{}l}{\emph{Captured lift at a quarter budget}} \\
% src: paper/data/run11-simulated-readers/simulated-readers-claude-opus-5-5-20261106-2dc461f52ea4.json ranking.n_units; paper/data/run11-simulated-readers/simulated-readers-claude-sonnet-5-5-20261106-2dc461f52ea4.json ranking.n_units; paper/data/run11-simulated-readers/simulated-readers-claude-haiku-5-5-20261106-2dc461f52ea4.json ranking.n_units
units ranked & 1,771 & 1,771 & 1,771 \\
% src: paper/data/run11-simulated-readers/simulated-readers-claude-opus-5-5-20261106-2dc461f52ea4.json ranking.panel.estimate; paper/data/run11-simulated-readers/simulated-readers-claude-sonnet-5-5-20261106-2dc461f52ea4.json ranking.panel.estimate; paper/data/run11-simulated-readers/simulated-readers-claude-haiku-5-5-20261106-2dc461f52ea4.json ranking.panel.estimate
plain panel & 0.333 & 0.288 & 0.306 \\
% src: paper/data/run11-simulated-readers/simulated-readers-claude-opus-5-5-20261106-2dc461f52ea4.json ranking.arm_count.estimate; paper/data/run11-simulated-readers/simulated-readers-claude-sonnet-5-5-20261106-2dc461f52ea4.json ranking.arm_count.estimate; paper/data/run11-simulated-readers/simulated-readers-claude-haiku-5-5-20261106-2dc461f52ea4.json ranking.arm_count.estimate
arm-count rule & 0.317 & 0.317 & 0.317 \\
% src: paper/data/run11-simulated-readers/simulated-readers-claude-opus-5-5-20261106-2dc461f52ea4.json ranking.matched_random.estimate; paper/data/run11-simulated-readers/simulated-readers-claude-sonnet-5-5-20261106-2dc461f52ea4.json ranking.matched_random.estimate; paper/data/run11-simulated-readers/simulated-readers-claude-haiku-5-5-20261106-2dc461f52ea4.json ranking.matched_random.estimate
random order matched on arm count & 0.274 & 0.271 & 0.276 \\
% src: paper/data/run11-simulated-readers/simulated-readers-claude-opus-5-5-20261106-2dc461f52ea4.json ranking.panel_minus_arm_count.estimate; paper/data/run11-simulated-readers/simulated-readers-claude-opus-5-5-20261106-2dc461f52ea4.json ranking.panel_minus_arm_count.lo; paper/data/run11-simulated-readers/simulated-readers-claude-opus-5-5-20261106-2dc461f52ea4.json ranking.panel_minus_arm_count.hi; paper/data/run11-simulated-readers/simulated-readers-claude-sonnet-5-5-20261106-2dc461f52ea4.json ranking.panel_minus_arm_count.estimate; paper/data/run11-simulated-readers/simulated-readers-claude-sonnet-5-5-20261106-2dc461f52ea4.json ranking.panel_minus_arm_count.lo; paper/data/run11-simulated-readers/simulated-readers-claude-sonnet-5-5-20261106-2dc461f52ea4.json ranking.panel_minus_arm_count.hi; paper/data/run11-simulated-readers/simulated-readers-claude-haiku-5-5-20261106-2dc461f52ea4.json ranking.panel_minus_arm_count.estimate; paper/data/run11-simulated-readers/simulated-readers-claude-haiku-5-5-20261106-2dc461f52ea4.json ranking.panel_minus_arm_count.lo; paper/data/run11-simulated-readers/simulated-readers-claude-haiku-5-5-20261106-2dc461f52ea4.json ranking.panel_minus_arm_count.hi
plain panel minus the arm-count rule & 0.0163 (-0.0168, 0.0472) & -0.0282 (-0.0647, 0.0042) & -0.0104 (-0.0475, 0.0134) \\
% src: paper/data/run11-simulated-readers/simulated-readers-claude-opus-5-5-20261106-2dc461f52ea4.json ranking.panel_minus_matched_random.estimate; paper/data/run11-simulated-readers/simulated-readers-claude-opus-5-5-20261106-2dc461f52ea4.json ranking.panel_minus_matched_random.lo; paper/data/run11-simulated-readers/simulated-readers-claude-opus-5-5-20261106-2dc461f52ea4.json ranking.panel_minus_matched_random.hi; paper/data/run11-simulated-readers/simulated-readers-claude-sonnet-5-5-20261106-2dc461f52ea4.json ranking.panel_minus_matched_random.estimate; paper/data/run11-simulated-readers/simulated-readers-claude-sonnet-5-5-20261106-2dc461f52ea4.json ranking.panel_minus_matched_random.lo; paper/data/run11-simulated-readers/simulated-readers-claude-sonnet-5-5-20261106-2dc461f52ea4.json ranking.panel_minus_matched_random.hi; paper/data/run11-simulated-readers/simulated-readers-claude-haiku-5-5-20261106-2dc461f52ea4.json ranking.panel_minus_matched_random.estimate; paper/data/run11-simulated-readers/simulated-readers-claude-haiku-5-5-20261106-2dc461f52ea4.json ranking.panel_minus_matched_random.lo; paper/data/run11-simulated-readers/simulated-readers-claude-haiku-5-5-20261106-2dc461f52ea4.json ranking.panel_minus_matched_random.hi
plain panel minus the matched random order & 0.0586 (0.0304, 0.0825) & 0.0180 (-0.0074, 0.0428) & 0.0301 (0.0024, 0.0480) \\
\midrule
\multicolumn{4}{@{}l}{\emph{Plain panels of two models give the same direction, share}} \\
% src: paper/data/run11-simulated-readers/simulated-readers-agreement-20261106-2dc461f52ea4.json pairs["claude-opus-5-5__claude-sonnet-5-5"].agree_share; paper/data/run11-simulated-readers/simulated-readers-agreement-20261106-2dc461f52ea4.json pairs["claude-opus-5-5__claude-haiku-5-5"].agree_share
with Opus &  & 0.712 & 0.649 \\
% src: paper/data/run11-simulated-readers/simulated-readers-agreement-20261106-2dc461f52ea4.json pairs["claude-sonnet-5-5__claude-haiku-5-5"].agree_share
with Sonnet &  &  & 0.654 \\
\midrule
\multicolumn{4}{@{}l}{\emph{Calls}} \\
% src: paper/data/run11-simulated-readers/simulated-readers-claude-opus-5-5-20261106-2dc461f52ea4.json counts.calls_planned; paper/data/run11-simulated-readers/simulated-readers-claude-sonnet-5-5-20261106-2dc461f52ea4.json counts.calls_planned; paper/data/run11-simulated-readers/simulated-readers-claude-haiku-5-5-20261106-2dc461f52ea4.json counts.calls_planned
planned & 9,007 & 9,007 & 9,007 \\
% src: paper/data/run11-simulated-readers/simulated-readers-claude-opus-5-5-20261106-2dc461f52ea4.json counts.calls_ok; paper/data/run11-simulated-readers/simulated-readers-claude-sonnet-5-5-20261106-2dc461f52ea4.json counts.calls_ok; paper/data/run11-simulated-readers/simulated-readers-claude-haiku-5-5-20261106-2dc461f52ea4.json counts.calls_ok
usable replies & 9,007 & 9,007 & 9,007 \\
% src: paper/data/run11-simulated-readers/simulated-readers-claude-opus-5-5-20261106-2dc461f52ea4.json counts.calls_invalid; paper/data/run11-simulated-readers/simulated-readers-claude-sonnet-5-5-20261106-2dc461f52ea4.json counts.calls_invalid; paper/data/run11-simulated-readers/simulated-readers-claude-haiku-5-5-20261106-2dc461f52ea4.json counts.calls_invalid
invalid after three attempts & 0 & 0 & 0 \\
% src: paper/data/run11-simulated-readers/simulated-readers-claude-opus-5-5-20261106-2dc461f52ea4.json counts.calls_refusal; paper/data/run11-simulated-readers/simulated-readers-claude-sonnet-5-5-20261106-2dc461f52ea4.json counts.calls_refusal; paper/data/run11-simulated-readers/simulated-readers-claude-haiku-5-5-20261106-2dc461f52ea4.json counts.calls_refusal
declined & 0 & 0 & 0 \\
\bottomrule
\end{tabular}

\end{table*}

\begin{table*}[t]\centering\footnotesize
% src: tables/fmt-readers-more-api.tex; paper/data/run12-simulated-readers-api/simulated-readers-claude-haiku-5-5-run12-api-20261106-2dc461f52ea4.json counts.calls_invalid 4, counts.calls_ok 9003; paper/data/run12-simulated-readers-api/simulated-readers-claude-opus-5-5-run12-api-20261106-2dc461f52ea4.json ranking.frac 0.25, ranking.resamples 500, plain.decisive.n 1634; paper/data/run9-visible-exposure/backtest-upworthy-20261106-2dc461f52ea4.json tables.aggregate[cell=headline_only,rung=always_reference].direction_accuracy_decisive (the row of the earliest-arm rule on decisive contrasts); run 11 and run 12 are the names of the two run folders
\caption{Measures of run 12 that carry no verdict: the rows of Table~\ref{tab:fmt-readers-more} for the same requests sent through the vendor's programming interface. Decisive contrasts are those whose observed effect exceeds 1.96 standard errors in size. On them the earliest-arm rule is itself right more often than on all contrasts, so the plain panel's row there is read against that rule and not against one half. The rule's row comes from the main run's report, on the same contrasts. The two halves of the panel are readers 1 to 5 and readers 6 to 10. The Brier skill of the direction map is for the direction of a contrast and is not comparable with the Brier skills of Section~\ref{sec:ladder-results}. Captured lift is the share of the total positive lift held by the one fourth of the units ranked first within each origin, where a unit is a test at one origin. Intervals are at 90 percent from 500 resamples of blocks and are not adjusted.}
\label{tab:fmt-readers-more-api}
% generated by paper/build_tables.py from the report keys named above each row (readers); it reads paper/data/run12-simulated-readers-api and, for one row, paper/data/run9-visible-exposure
\begin{tabular}{p{6.6cm}rrr}
\toprule
measure & Opus & Sonnet & Haiku \\
\midrule
% src: paper/data/run12-simulated-readers-api/simulated-readers-claude-opus-5-5-run12-api-20261106-2dc461f52ea4.json plain.accuracy_minus_half.estimate; paper/data/run12-simulated-readers-api/simulated-readers-claude-opus-5-5-run12-api-20261106-2dc461f52ea4.json plain.accuracy_minus_half.lo; paper/data/run12-simulated-readers-api/simulated-readers-claude-opus-5-5-run12-api-20261106-2dc461f52ea4.json plain.accuracy_minus_half.hi; paper/data/run12-simulated-readers-api/simulated-readers-claude-sonnet-5-5-run12-api-20261106-2dc461f52ea4.json plain.accuracy_minus_half.estimate; paper/data/run12-simulated-readers-api/simulated-readers-claude-sonnet-5-5-run12-api-20261106-2dc461f52ea4.json plain.accuracy_minus_half.lo; paper/data/run12-simulated-readers-api/simulated-readers-claude-sonnet-5-5-run12-api-20261106-2dc461f52ea4.json plain.accuracy_minus_half.hi; paper/data/run12-simulated-readers-api/simulated-readers-claude-haiku-5-5-run12-api-20261106-2dc461f52ea4.json plain.accuracy_minus_half.estimate; paper/data/run12-simulated-readers-api/simulated-readers-claude-haiku-5-5-run12-api-20261106-2dc461f52ea4.json plain.accuracy_minus_half.lo; paper/data/run12-simulated-readers-api/simulated-readers-claude-haiku-5-5-run12-api-20261106-2dc461f52ea4.json plain.accuracy_minus_half.hi
plain panel minus one half & 0.0742 (0.0623, 0.0867) & 0.0279 (0.0176, 0.0404) & 0.0223 (0.0103, 0.0344) \\
% src: paper/data/run12-simulated-readers-api/simulated-readers-claude-opus-5-5-run12-api-20261106-2dc461f52ea4.json plain.decisive.n; paper/data/run12-simulated-readers-api/simulated-readers-claude-sonnet-5-5-run12-api-20261106-2dc461f52ea4.json plain.decisive.n; paper/data/run12-simulated-readers-api/simulated-readers-claude-haiku-5-5-run12-api-20261106-2dc461f52ea4.json plain.decisive.n
decisive contrasts & 1,634 & 1,634 & 1,634 \\
% src: paper/data/run12-simulated-readers-api/simulated-readers-claude-opus-5-5-run12-api-20261106-2dc461f52ea4.json plain.decisive.accuracy; paper/data/run12-simulated-readers-api/simulated-readers-claude-sonnet-5-5-run12-api-20261106-2dc461f52ea4.json plain.decisive.accuracy; paper/data/run12-simulated-readers-api/simulated-readers-claude-haiku-5-5-run12-api-20261106-2dc461f52ea4.json plain.decisive.accuracy
plain panel, share matched on decisive contrasts & 0.6506 & 0.5863 & 0.5624 \\
% src: paper/data/run9-visible-exposure/backtest-upworthy-20261106-2dc461f52ea4.json tables.aggregate[cell=headline_only,rung=always_reference].direction_accuracy_decisive; paper/data/run9-visible-exposure/backtest-upworthy-20261106-2dc461f52ea4.json tables.aggregate[cell=headline_only,rung=always_reference].direction_accuracy_decisive; paper/data/run9-visible-exposure/backtest-upworthy-20261106-2dc461f52ea4.json tables.aggregate[cell=headline_only,rung=always_reference].direction_accuracy_decisive
the earliest-arm rule, share matched on decisive contrasts & 0.6481 & 0.6481 & 0.6481 \\
% src: paper/data/run12-simulated-readers-api/simulated-readers-claude-opus-5-5-run12-api-20261106-2dc461f52ea4.json plain.split_half.agree_share; paper/data/run12-simulated-readers-api/simulated-readers-claude-sonnet-5-5-run12-api-20261106-2dc461f52ea4.json plain.split_half.agree_share; paper/data/run12-simulated-readers-api/simulated-readers-claude-haiku-5-5-run12-api-20261106-2dc461f52ea4.json plain.split_half.agree_share
the two halves of the panel give the same direction, share & 0.747 & 0.763 & 0.669 \\
% src: paper/data/run12-simulated-readers-api/simulated-readers-claude-opus-5-5-run12-api-20261106-2dc461f52ea4.json plain.ties; paper/data/run12-simulated-readers-api/simulated-readers-claude-sonnet-5-5-run12-api-20261106-2dc461f52ea4.json plain.ties; paper/data/run12-simulated-readers-api/simulated-readers-claude-haiku-5-5-run12-api-20261106-2dc461f52ea4.json plain.ties
contrasts on which the plain panel is tied & 1 & 1 & 2 \\
% src: paper/data/run12-simulated-readers-api/simulated-readers-claude-opus-5-5-run12-api-20261106-2dc461f52ea4.json plain.no_reader; paper/data/run12-simulated-readers-api/simulated-readers-claude-sonnet-5-5-run12-api-20261106-2dc461f52ea4.json plain.no_reader; paper/data/run12-simulated-readers-api/simulated-readers-claude-haiku-5-5-run12-api-20261106-2dc461f52ea4.json plain.no_reader
contrasts left with no reader & 0 & 0 & 0 \\
% src: paper/data/run12-simulated-readers-api/simulated-readers-claude-opus-5-5-run12-api-20261106-2dc461f52ea4.json methods.direction_map.brier_skill; paper/data/run12-simulated-readers-api/simulated-readers-claude-sonnet-5-5-run12-api-20261106-2dc461f52ea4.json methods.direction_map.brier_skill; paper/data/run12-simulated-readers-api/simulated-readers-claude-haiku-5-5-run12-api-20261106-2dc461f52ea4.json methods.direction_map.brier_skill
direction map, Brier skill & 0.0208 & 0.0035 & -0.0002 \\
% src: paper/data/run12-simulated-readers-api/simulated-readers-claude-opus-5-5-run12-api-20261106-2dc461f52ea4.json baselines.text_lr_minus_always_reference.estimate; paper/data/run12-simulated-readers-api/simulated-readers-claude-opus-5-5-run12-api-20261106-2dc461f52ea4.json baselines.text_lr_minus_always_reference.lo; paper/data/run12-simulated-readers-api/simulated-readers-claude-opus-5-5-run12-api-20261106-2dc461f52ea4.json baselines.text_lr_minus_always_reference.hi; paper/data/run12-simulated-readers-api/simulated-readers-claude-sonnet-5-5-run12-api-20261106-2dc461f52ea4.json baselines.text_lr_minus_always_reference.estimate; paper/data/run12-simulated-readers-api/simulated-readers-claude-sonnet-5-5-run12-api-20261106-2dc461f52ea4.json baselines.text_lr_minus_always_reference.lo; paper/data/run12-simulated-readers-api/simulated-readers-claude-sonnet-5-5-run12-api-20261106-2dc461f52ea4.json baselines.text_lr_minus_always_reference.hi; paper/data/run12-simulated-readers-api/simulated-readers-claude-haiku-5-5-run12-api-20261106-2dc461f52ea4.json baselines.text_lr_minus_always_reference.estimate; paper/data/run12-simulated-readers-api/simulated-readers-claude-haiku-5-5-run12-api-20261106-2dc461f52ea4.json baselines.text_lr_minus_always_reference.lo; paper/data/run12-simulated-readers-api/simulated-readers-claude-haiku-5-5-run12-api-20261106-2dc461f52ea4.json baselines.text_lr_minus_always_reference.hi
text classifier minus the earliest-arm rule & -0.0216 (-0.0365, -0.0070) & -0.0216 (-0.0365, -0.0070) & -0.0216 (-0.0365, -0.0070) \\
\midrule
\multicolumn{4}{@{}l}{\emph{Captured lift at a quarter budget}} \\
% src: paper/data/run12-simulated-readers-api/simulated-readers-claude-opus-5-5-run12-api-20261106-2dc461f52ea4.json ranking.n_units; paper/data/run12-simulated-readers-api/simulated-readers-claude-sonnet-5-5-run12-api-20261106-2dc461f52ea4.json ranking.n_units; paper/data/run12-simulated-readers-api/simulated-readers-claude-haiku-5-5-run12-api-20261106-2dc461f52ea4.json ranking.n_units
units ranked & 1,771 & 1,771 & 1,771 \\
% src: paper/data/run12-simulated-readers-api/simulated-readers-claude-opus-5-5-run12-api-20261106-2dc461f52ea4.json ranking.panel.estimate; paper/data/run12-simulated-readers-api/simulated-readers-claude-sonnet-5-5-run12-api-20261106-2dc461f52ea4.json ranking.panel.estimate; paper/data/run12-simulated-readers-api/simulated-readers-claude-haiku-5-5-run12-api-20261106-2dc461f52ea4.json ranking.panel.estimate
plain panel & 0.339 & 0.315 & 0.321 \\
% src: paper/data/run12-simulated-readers-api/simulated-readers-claude-opus-5-5-run12-api-20261106-2dc461f52ea4.json ranking.arm_count.estimate; paper/data/run12-simulated-readers-api/simulated-readers-claude-sonnet-5-5-run12-api-20261106-2dc461f52ea4.json ranking.arm_count.estimate; paper/data/run12-simulated-readers-api/simulated-readers-claude-haiku-5-5-run12-api-20261106-2dc461f52ea4.json ranking.arm_count.estimate
arm-count rule & 0.317 & 0.317 & 0.317 \\
% src: paper/data/run12-simulated-readers-api/simulated-readers-claude-opus-5-5-run12-api-20261106-2dc461f52ea4.json ranking.matched_random.estimate; paper/data/run12-simulated-readers-api/simulated-readers-claude-sonnet-5-5-run12-api-20261106-2dc461f52ea4.json ranking.matched_random.estimate; paper/data/run12-simulated-readers-api/simulated-readers-claude-haiku-5-5-run12-api-20261106-2dc461f52ea4.json ranking.matched_random.estimate
random order matched on arm count & 0.277 & 0.273 & 0.275 \\
% src: paper/data/run12-simulated-readers-api/simulated-readers-claude-opus-5-5-run12-api-20261106-2dc461f52ea4.json ranking.panel_minus_arm_count.estimate; paper/data/run12-simulated-readers-api/simulated-readers-claude-opus-5-5-run12-api-20261106-2dc461f52ea4.json ranking.panel_minus_arm_count.lo; paper/data/run12-simulated-readers-api/simulated-readers-claude-opus-5-5-run12-api-20261106-2dc461f52ea4.json ranking.panel_minus_arm_count.hi; paper/data/run12-simulated-readers-api/simulated-readers-claude-sonnet-5-5-run12-api-20261106-2dc461f52ea4.json ranking.panel_minus_arm_count.estimate; paper/data/run12-simulated-readers-api/simulated-readers-claude-sonnet-5-5-run12-api-20261106-2dc461f52ea4.json ranking.panel_minus_arm_count.lo; paper/data/run12-simulated-readers-api/simulated-readers-claude-sonnet-5-5-run12-api-20261106-2dc461f52ea4.json ranking.panel_minus_arm_count.hi; paper/data/run12-simulated-readers-api/simulated-readers-claude-haiku-5-5-run12-api-20261106-2dc461f52ea4.json ranking.panel_minus_arm_count.estimate; paper/data/run12-simulated-readers-api/simulated-readers-claude-haiku-5-5-run12-api-20261106-2dc461f52ea4.json ranking.panel_minus_arm_count.lo; paper/data/run12-simulated-readers-api/simulated-readers-claude-haiku-5-5-run12-api-20261106-2dc461f52ea4.json ranking.panel_minus_arm_count.hi
plain panel minus the arm-count rule & 0.0227 (-0.0065, 0.0543) & -0.0014 (-0.0378, 0.0272) & 0.0040 (-0.0328, 0.0365) \\
% src: paper/data/run12-simulated-readers-api/simulated-readers-claude-opus-5-5-run12-api-20261106-2dc461f52ea4.json ranking.panel_minus_matched_random.estimate; paper/data/run12-simulated-readers-api/simulated-readers-claude-opus-5-5-run12-api-20261106-2dc461f52ea4.json ranking.panel_minus_matched_random.lo; paper/data/run12-simulated-readers-api/simulated-readers-claude-opus-5-5-run12-api-20261106-2dc461f52ea4.json ranking.panel_minus_matched_random.hi; paper/data/run12-simulated-readers-api/simulated-readers-claude-sonnet-5-5-run12-api-20261106-2dc461f52ea4.json ranking.panel_minus_matched_random.estimate; paper/data/run12-simulated-readers-api/simulated-readers-claude-sonnet-5-5-run12-api-20261106-2dc461f52ea4.json ranking.panel_minus_matched_random.lo; paper/data/run12-simulated-readers-api/simulated-readers-claude-sonnet-5-5-run12-api-20261106-2dc461f52ea4.json ranking.panel_minus_matched_random.hi; paper/data/run12-simulated-readers-api/simulated-readers-claude-haiku-5-5-run12-api-20261106-2dc461f52ea4.json ranking.panel_minus_matched_random.estimate; paper/data/run12-simulated-readers-api/simulated-readers-claude-haiku-5-5-run12-api-20261106-2dc461f52ea4.json ranking.panel_minus_matched_random.lo; paper/data/run12-simulated-readers-api/simulated-readers-claude-haiku-5-5-run12-api-20261106-2dc461f52ea4.json ranking.panel_minus_matched_random.hi
plain panel minus the matched random order & 0.0621 (0.0364, 0.0844) & 0.0418 (0.0114, 0.0621) & 0.0455 (0.0152, 0.0686) \\
\midrule
\multicolumn{4}{@{}l}{\emph{Plain panels of two models give the same direction, share}} \\
% src: paper/data/run12-simulated-readers-api/simulated-readers-agreement-run12-api-20261106-2dc461f52ea4.json pairs["claude-opus-5-5__claude-sonnet-5-5"].agree_share; paper/data/run12-simulated-readers-api/simulated-readers-agreement-run12-api-20261106-2dc461f52ea4.json pairs["claude-opus-5-5__claude-haiku-5-5"].agree_share
with Opus &  & 0.717 & 0.598 \\
% src: paper/data/run12-simulated-readers-api/simulated-readers-agreement-run12-api-20261106-2dc461f52ea4.json pairs["claude-sonnet-5-5__claude-haiku-5-5"].agree_share
with Sonnet &  &  & 0.626 \\
\midrule
\multicolumn{4}{@{}l}{\emph{Calls}} \\
% src: paper/data/run12-simulated-readers-api/simulated-readers-claude-opus-5-5-run12-api-20261106-2dc461f52ea4.json counts.calls_planned; paper/data/run12-simulated-readers-api/simulated-readers-claude-sonnet-5-5-run12-api-20261106-2dc461f52ea4.json counts.calls_planned; paper/data/run12-simulated-readers-api/simulated-readers-claude-haiku-5-5-run12-api-20261106-2dc461f52ea4.json counts.calls_planned
planned & 9,007 & 9,007 & 9,007 \\
% src: paper/data/run12-simulated-readers-api/simulated-readers-claude-opus-5-5-run12-api-20261106-2dc461f52ea4.json counts.calls_ok; paper/data/run12-simulated-readers-api/simulated-readers-claude-sonnet-5-5-run12-api-20261106-2dc461f52ea4.json counts.calls_ok; paper/data/run12-simulated-readers-api/simulated-readers-claude-haiku-5-5-run12-api-20261106-2dc461f52ea4.json counts.calls_ok
usable replies & 9,007 & 9,007 & 9,003 \\
% src: paper/data/run12-simulated-readers-api/simulated-readers-claude-opus-5-5-run12-api-20261106-2dc461f52ea4.json counts.calls_invalid; paper/data/run12-simulated-readers-api/simulated-readers-claude-sonnet-5-5-run12-api-20261106-2dc461f52ea4.json counts.calls_invalid; paper/data/run12-simulated-readers-api/simulated-readers-claude-haiku-5-5-run12-api-20261106-2dc461f52ea4.json counts.calls_invalid
invalid after three attempts & 0 & 0 & 4 \\
% src: paper/data/run12-simulated-readers-api/simulated-readers-claude-opus-5-5-run12-api-20261106-2dc461f52ea4.json counts.calls_refusal; paper/data/run12-simulated-readers-api/simulated-readers-claude-sonnet-5-5-run12-api-20261106-2dc461f52ea4.json counts.calls_refusal; paper/data/run12-simulated-readers-api/simulated-readers-claude-haiku-5-5-run12-api-20261106-2dc461f52ea4.json counts.calls_refusal
declined & 0 & 0 & 0 \\
\bottomrule
\end{tabular}

\end{table*}

% src: tables/fmt-readers-more.tex rows plain panel minus one half: Opus 0.0742 (0.0616, 0.0878); plain panel minus the arm-count rule: Opus 0.0163 (-0.0168, 0.0472), Sonnet -0.0282 (-0.0647, 0.0042), Haiku -0.0104 (-0.0475, 0.0134); text classifier minus the earliest-arm rule -0.0216 (-0.0365, -0.0070); with Opus 0.712 and 0.649; tables/fmt-readers-more-api.tex plain panel minus the arm-count rule: Opus 0.0227 (-0.0065, 0.0543), Sonnet -0.0014 (-0.0378, 0.0272), Haiku 0.0040 (-0.0328, 0.0365); rows plain panel, share matched on decisive contrasts and the earliest-arm rule, share matched on decisive contrasts, in both tables (Opus below the rule in run 11 and above it in run 12, Sonnet and Haiku below it in both; no report holds an interval for these differences; in words only); run 11 and run 12 are the names of the two run folders
\paragraph{Reading the remaining measures.} Five points of Tables~\ref{tab:fmt-readers-more} and~\ref{tab:fmt-readers-more-api} bear on the body's readings. First, every plain panel is above one half, by 0.074 for the primary model, and that is a low bar: the earliest-arm rule clears it without reading a headline. Second, on these rows the text classifier is itself below the earliest-arm rule, by 0.022, with an interval that lies below zero. Third, as a ranking of tests the plain panel is not shown to capture more lift than counting a test's arms: for each model and in each run the interval of that difference contains zero. Fourth, the panels of two models give the same direction on 0.649 to 0.712 of the contrasts in run 11, which is less than the two halves of one panel of the primary model. Fifth, on the decisive contrasts the plain panel of the primary model is close to the earliest-arm rule, below it in run 11 and above it in run 12, and the panels of the two smaller models are below it in both runs. No interval was computed for these differences, so they describe and carry no verdict.
